%% 4-13 空間の座標と3次元のベクトル（第13回：組 V4-7。ページ制限なし。偶数ページ）
%% ほかの巻から指されている形：
%%   3-2「空間の中の場所は、高さの z を加えた3つの数 (x, y, z) で表します（4-13）」
%%   第4巻の約束「工学の専門書では太字 a で書くことが多い。物理シリーズはこの本と同じ矢印」をこの章で紹介する（第13回の決定）
\makeatletter
% 斜めから見た空間の図（x 軸を左下へ、y 軸を右へ、z 軸を上へかく）  [ivn3d=1目もりの長さ]
\tikzset{ivn3d/.style={x={({-0.49*#1},{-0.34*#1})}, y={(#1,0)}, z={(0,#1)}}}
% 式を、まん中に1行で置く  \bkivnline{…}
\newcommand\bkivnline[1]{\par\addvspace{3pt}{\centering #1\par}\addvspace{3pt}\noindent\ignorespaces}
% 説明の段の見出し（丸数字＋太字の1行）  \bkivnstep{番号}{見出し}
\newcommand\bkivnstep[2]{\par\addvspace{6pt}\Needspace{5\baselineskip}\noindent
  {\bkcircled[bkOrange]{#1}\enspace\sffamily\bfseries\color{bkOrange!75!black}#2}\par\nobreak\vspace{2pt}}
% 小見出し（それだけで1行）
\newcommand\bkivnhead[1]{\par\addvspace{4pt}\Needspace{6\baselineskip}\noindent\textbf{#1}\par\nobreak}
% 3本の軸  \bkivnaxes{x の長さ}{y の長さ}{z の長さ}
\newcommand\bkivnaxes[3]{%
  \draw[bkGray, line width=0.6pt, -{Stealth[length=4pt]}] (0,0,0) -- (#1,0,0) node[anchor=east, inner sep=2pt, font=\sffamily\small, text=bkInk] {$x$};
  \draw[bkGray, line width=0.6pt, -{Stealth[length=4pt]}] (0,0,0) -- (0,#2,0) node[anchor=west, inner sep=2pt, font=\sffamily\small, text=bkInk] {$y$};
  \draw[bkGray, line width=0.6pt, -{Stealth[length=4pt]}] (0,0,0) -- (0,0,#3) node[anchor=south, inner sep=2pt, font=\sffamily\small, text=bkInk] {$z$};}
% 原点から (a, b, c) までの直方体の、見えている辺（実線）と、原点から出る辺（点線）
\newcommand\bkivnbox[4]{%
  \draw[#4, line width=0.7pt] (#1,0,0) -- (#1,#2,0) -- (0,#2,0) -- (0,#2,#3) -- (0,0,#3) -- (#1,0,#3) -- cycle;
  \draw[#4, line width=0.7pt] (#1,0,#3) -- (#1,#2,#3) -- (0,#2,#3) (#1,#2,#3) -- (#1,#2,0);
  \draw[#4, line width=0.5pt, dashed] (0,0,0) -- (#1,0,0) (0,0,0) -- (0,#2,0) (0,0,0) -- (0,0,#3);}
% 小さなドローン（canvas の mm でかく）  \bkivndrone{(点)}
\newcommand\bkivndrone[1]{\begin{scope}[shift={#1}, scale=1.35]
  \draw[bkInk!70, line width=0.7pt] (-4mm,1.4mm) -- (4mm,1.4mm);
  \draw[bkInk!70, line width=0.6pt] (-4mm,1.4mm) -- (-4mm,2.1mm) (4mm,1.4mm) -- (4mm,2.1mm);
  \fill[bkGray!35] (-4mm,2.3mm) ellipse (2.4mm and 0.55mm);
  \fill[bkGray!35] (4mm,2.3mm) ellipse (2.4mm and 0.55mm);
  \filldraw[fill=bkNavy!75, draw=bkNavy, line width=0.4pt, rounded corners=0.8mm] (-2.6mm,0.4mm) rectangle (2.6mm,2.2mm);
  \fill[bkYellow] (0,1.3mm) circle (0.45mm);
\end{scope}}
% 右手（3本の指を、たがいに直角に開いた形）。canvas の mm でかく  \bkivnhand
\newcommand\bkivnfinger[3]{% #1 から #2 へ、太さ #3
  \draw[bkOrange!70!black, line width=#3, line cap=round] #1 -- #2;
  \draw[bkOrange!22, line width=#3-1pt, line cap=round] #1 -- #2;}
\makeatother

\bkchapter{4}{13}{空間の座標と\\3次元のベクトル}{\bklevel{4}}{空間の座標と3次元のベクトル}
\bkmichishirube{3-2,4-4,4-11,4-12}{4-14,4-15,7-7,8-5}

\begin{goal}
\begin{itemize}[label={\color{\bkTT}●}, leftmargin=1.6\zw, itemsep=0pt, topsep=1pt]
\item \bklevel{3}\enspace 空間の点を3つの数 $(x,\ y,\ z)$ で表し、2点の距離を求められる。
\item \bklevel{4}\enspace 3次元のベクトルを成分で計算し、内積からなす角を求められる。
\item \bklevel{4}\enspace 基本ベクトル・単位ベクトル・法線ベクトルで、向きや平面を表せる。
\end{itemize}
\end{goal}

\begin{hatena}
\noindent\begin{minipage}[c]{0.56\linewidth}
部屋の中を、小さなドローンが飛んでいます。となりの部屋にいる友だちに、ドローンがいまどこにいるかを、電話で正しく伝えたい。

床の上の場所なら、3-2の座標のように「右へいくつ、奥へいくつ」の2つの数で言えました。でも、ドローンは宙にういています。

ドローンの場所を言うには、数がいくつあればいいかな？
\end{minipage}\hfill
\begin{minipage}[c]{0.4\linewidth}
\centering
\begin{tikzpicture}[ivn3d=6.4mm]
  \fill[bkBlue!6] (0,0,0) -- (0,5,0) -- (0,5,2.8) -- (0,0,2.8) -- cycle;
  \fill[bkGreen!7] (0,0,0) -- (4,0,0) -- (4,0,2.8) -- (0,0,2.8) -- cycle;
  \fill[bkYellow!14] (0,0,0) -- (4,0,0) -- (4,5,0) -- (0,5,0) -- cycle;
  \draw[bkGray!80, line width=0.6pt] (4,0,0) -- (4,5,0) -- (0,5,0) (0,5,0) -- (0,5,2.8) -- (0,0,2.8) -- (4,0,2.8) -- (4,0,0) (0,0,0) -- (0,0,2.8);
  \draw[bkGray!80, line width=0.6pt] (0,0,0) -- (4,0,0) (0,0,0) -- (0,5,0);
  \draw[bkInk!60, line width=0.5pt, dashed] (2,3,1.6) -- (2,3,0);
  \fill[bkInk!35] (2,3,0) ellipse (2.2mm and 0.7mm);
  \bkivndrone{(2,3,1.6)}
  \node[font=\sffamily\footnotesize, text=bkInk, anchor=south west, inner sep=1pt] at ($(4,0,0)+(1mm,0.5mm)$) {床};
  \node[font=\sffamily\footnotesize, text=bkInk, anchor=west] at ($(2,3,1.6)+(5mm,1.5mm)$) {?};
\end{tikzpicture}
\end{minipage}
\end{hatena}

\noindent 2-2の数直線では、1つの数で場所が決まりました。3-2の座標では、2つの数で平面の上の場所が決まりました。この章では、部屋や空の中のような「空間」の場所を数で表します。そして、4-10〜4-12で学んだベクトル（向きと大きさをもつ量）を、空間でも使えるようにします。

%% ---------------- やってみよう・思い出そう ----------------
\begin{yattemiyou}[部屋の中の物の場所を、数ではかろう]
\noindent\begin{minipage}[c]{0.6\linewidth}
\textbf{用意するもの}\par
メジャー（2 m 以上）、ふせん、ノート
\begin{enumerate}[label={\bkcircled[bkGreen]{\arabic*}}, leftmargin=2.4\zw, itemsep=1pt, topsep=3pt]
\item 部屋の床の1つのすみを「原点O」に決めよう。
\item すみから出る3本の線（床とかべの境目2本と、かべとかべの境目）にふせんで名前をつけよう。手前へ $x$、右へ $y$、上へ $z$。
\item 机の角、電灯、かべの時計の場所を、$x$、$y$、$z$ の順にメジャーではかろう。高さ $z$ は、物の真下の床に印をつけてからはかるとよい。
\end{enumerate}
\end{minipage}\hfill
\begin{minipage}[c]{0.36\linewidth}
\centering
\begin{tikzpicture}[ivn3d=5mm]
  \bkivnaxes{4.6}{5.4}{3.6}
  \fill[bkInk] (0,0,0) circle (0.6mm);
  \node[font=\sffamily\footnotesize, anchor=east, inner sep=2pt] at (0,0,0) {O};
  \draw[bkGray!70, line width=0.5pt] (4,0,0) -- (4,5,0) -- (0,5,0);
  % 机（上の面の角を (1.2, 0.8, 0.7) とする。机は (1.2〜2.4, 0.8〜2.4) 、高さ 0.7）
  \filldraw[fill=bkOrange!25, draw=bkOrange!80!black, line width=0.5pt] (2.4,0.8,0.7) -- (2.4,2.4,0.7) -- (1.2,2.4,0.7) -- (1.2,0.8,0.7) -- cycle;
  \draw[bkOrange!80!black, line width=0.8pt] (2.4,0.8,0.7) -- (2.4,0.8,0) (2.4,2.4,0.7) -- (2.4,2.4,0);
  \fill[bkRed] (1.2,0.8,0.7) circle (0.6mm);
  % 電灯
  \draw[bkGray, line width=0.5pt] (2,1.8,2.9) -- (2,1.8,2.5);
  \filldraw[fill=bkYellow, draw=bkYellow!50!black, line width=0.4pt] ($(2,1.8,2.5)+(-2mm,0)$) -- ($(2,1.8,2.5)+(2mm,0)$) -- ($(2,1.8,2.5)+(1mm,1.5mm)$) -- ($(2,1.8,2.5)+(-1mm,1.5mm)$) -- cycle;
  % 時計（y=0 のかべ）
  \filldraw[fill=white, draw=bkInk, line width=0.5pt] (1.5,0,2) circle (1.6mm);
  \draw[bkInk, line width=0.4pt] (1.5,0,2) -- ++(0,0,0.2) (1.5,0,2) -- ++(0.25mm,0.9mm);
\end{tikzpicture}
\end{minipage}
\par\vspace{8pt}\noindent ソラさんの部屋ではかった結果です（単位は m）。
\begin{bkcenter}
\begin{bktable}{colspec={Q[l,m,wd=30mm] *{3}{Q[c,m,wd=20mm]}}, rowsep=1.5pt, row{1}={font=\sffamily\small}}
 & 手前へ $x$ & 右へ $y$ & 上へ $z$ \\
机の角（赤い点） & 1.2 & 0.8 & 0.7 \\
電灯 & 2.0 & 1.8 & 2.5 \\
かべの時計 & 1.5 & 0 & 2.0 \\
\end{bktable}
\end{bkcenter}
\end{yattemiyou}

\noindent どの物も、3つの数がそろうと、部屋の中の場所が1つに決まりました。$x$ と $y$ だけでは、「床のどこの真上か」しかわかりません。机の角と、その真上の天井の点は、$x$ と $y$ が同じで、$z$ だけがちがいます。かべの時計は、$y=0$ のかべにかかっているので、$y$ が0です。

\begin{omoidasou}
\bkrecall{3-2}{平面の点は、原点から右へ $x$、上へ $y$ の所で、$(x,\ y)$ と書く。原点と＋の向きは自分で決め、決めたら変えない。}
\bkrecall{4-4}{三平方の定理：直角三角形で（斜辺）$^2$＝（たて）$^2$＋（横）$^2$。\\2点の距離は $d=\sqrt{(x_2-x_1)^2+(y_2-y_1)^2}$。}
\bkrecall{4-11}{ベクトルは成分 $(a_x,\ a_y)$ で書ける。大きさは $\sqrt{a_x^2+a_y^2}$。和・差・実数倍は成分ごとに計算する。大きさが1のベクトルが単位ベクトル。$x$ の向きの単位ベクトルは $\vec e_x=(1,\ 0)$、$y$ の向きは $\vec e_y=(0,\ 1)$。}
\bkrecall{4-12}{内積 $\vec a\cdot\vec b=|\vec a||\vec b|\cos\theta=a_xb_x+a_yb_y$。直交すると0。仕事は $W=\vec F\cdot\vec s$。}
\end{omoidasou}

\noindent この章では、角度や長さの計算で $\approx$ を使います。$\approx$ は「ほぼ等しい」というしるしです（5-8でくわしく）。

%% ---------------- 数学の道具 ----------------
\bksection{doguu}{数学の道具}

\bksubsection{1}{空間の場所は、3つの数で決まる}

\leavevmode\label{bk:4-13:zahyo}場所を言うのに、数がいくつ必要かを、3つの場面でくらべてみます。

\begin{bkfigure}
\begin{tikzpicture}[x=1mm, y=1mm]
  % 1つの数：まっすぐな廊下
  \node[font=\sffamily\small\bfseries, anchor=south] at (19,32) {まっすぐな廊下};
  \fill[bkGray!12] (0,15) rectangle (38,23);
  \draw[bkInk, line width=0.8pt, -{Stealth[length=4pt]}] (0,19) -- (40,19);
  \foreach \k in {0,...,7} {\draw[bkInk, line width=0.5pt] (\k*5,17.8) -- (\k*5,20.2);}
  \node[font=\sffamily\footnotesize, anchor=north] at (0,14.4) {0};
  \node[font=\sffamily\footnotesize, anchor=north] at (25,14.4) {5};
  \fill[bkRed] (25,19) circle (1.2);
  \node[font=\sffamily\small, align=center] at (19,2) {入り口から 5 m\\[1pt]数は\textbf{1つ}};
  % 2つの数：床（平面）
  \begin{scope}[shift={(48,0)}]
  \node[font=\sffamily\small\bfseries, anchor=south] at (19,32) {教室の床};
  \draw[bkLine, line width=0.3pt] (0,8) grid[step=5] (35,28);
  \draw[bkInk, line width=0.7pt, -{Stealth[length=4pt]}] (0,8) -- (38,8);
  \draw[bkInk, line width=0.7pt, -{Stealth[length=4pt]}] (0,8) -- (0,30);
  \fill[bkRed] (20,23) circle (1.2);
  \draw[bkRed!80!black, line width=0.5pt, dashed] (20,8) -- (20,23) -- (0,23);
  \node[font=\sffamily\footnotesize, anchor=south west, inner sep=1pt] at (20.8,23.6) {$(4,\ 3)$};
  \node[font=\sffamily\small, align=center] at (19,2) {右へ4、奥へ3\\[1pt]数は\textbf{2つ}};
  \end{scope}
  % 3つの数：部屋（空間）
  \begin{scope}[shift={(100,13)}, ivn3d=5mm]
  \node[font=\sffamily\small\bfseries, anchor=south] at ($(0,2,0)+(0,19mm)$) {部屋の中};
  \bkivnaxes{2.8}{5.4}{3.4}
  \bkivnbox{3}{4}{2}{bkGray}
  \fill[bkRed] (3,4,2) circle (1.2mm);
  \node[font=\sffamily\footnotesize, anchor=south, inner sep=1.5pt, fill=white] at ($(3,4,2)+(0,1mm)$) {$(3,\ 4,\ 2)$};
  \node[font=\sffamily\small, align=center] at ($(0,2.4,0)+(0,-11mm)$) {手前へ3、右へ4、上へ2\\[1pt]数は\textbf{3つ}};
  \end{scope}
\end{tikzpicture}
\bkfigcaption{動ける向きが1つ（前・後ろ）なら数は1つ、2つ（前後・左右）なら2つ、3つ（前後・左右・上下）なら3つ。}\label{bk:4-13:fig-123}
\end{bkfigure}

まっすぐな廊下では、前か後ろにしか動けません。だから「入り口から何m」の1つの数で場所が決まります。床の上では、前後と左右の2つの向きに動けるので、数が2つ必要です（3-2）。部屋の中を飛ぶドローンは、前後・左右に加えて、上下にも動けます。だから数が3つ必要です。

\begin{bkcenter}
\begin{bktable}{colspec={Q[l,m,wd=26mm] Q[c,m,wd=34mm] Q[c,m,wd=22mm] Q[c,m,wd=30mm]}, rowsep=1.5pt, row{1}={font=\sffamily\small}}
 & 動ける向き & 必要な数 & 名前 \\
直線（廊下） & 前後 & 1つ & \textbf{1次元} \\
平面（床） & 前後・左右 & 2つ & \textbf{2次元} \\
空間（部屋） & 前後・左右・上下 & 3つ & \textbf{3次元} \\
\end{bktable}
\bktabcaption{場所を決めるのに必要な数の個数を、次元（じげん）という。}
\end{bkcenter}

場所を決めるのに必要な数の個数を\textbf{次元}（じげん）といいます。\bkidx{じげん}{次元}直線は1次元、平面は2次元、わたしたちのいる\textbf{空間}（くうかん）は3次元です。映画の「3D」の D は、次元を表す英語 dimension の頭文字です。

\bkivnstep{1}{3本の軸と、3つの数}
平面では、原点Oで直角に交わる2本の数直線（$x$ 軸と $y$ 軸）を使いました（3-2）。空間では、原点Oで\textbf{たがいに直角に交わる3本の数直線}を使います。3本めを $z$ 軸といい、ふつうは上向きにとります。

\noindent\begin{minipage}[c]{0.5\linewidth}
\centering
\begin{tikzpicture}[ivn3d=7.5mm]
  \bkivnaxes{4.6}{5.4}{3.3}
  \bkivnbox{3}{4}{2}{bkGray!70}
  \draw[bkRed, line width=1.6pt, -{Stealth[length=5pt]}] (0,0,0) -- (3,0,0);
  \draw[bkBlue, line width=1.6pt, -{Stealth[length=5pt]}] (3,0,0) -- (3,4,0);
  \draw[bkGreen!70!black, line width=1.6pt, -{Stealth[length=5pt]}] (3,4,0) -- (3,4,2);
  \fill[bkInk] (0,0,0) circle (0.7mm);
  \node[font=\sffamily\footnotesize, anchor=south east, inner sep=1.5pt] at (0,0,0) {O};
  \fill[bkInk] (3,4,2) circle (0.9mm);
  \node[font=\sffamily\small, anchor=south west, inner sep=1.5pt] at (3,4,2) {P};
  \node[font=\sffamily\small, text=bkRed!80!black, anchor=north west, inner sep=1pt] at ($(1.5,0,0)+(0.8mm,-0.8mm)$) {①};
  \node[font=\sffamily\small, text=bkBlue!80!black, anchor=north, inner sep=1pt] at ($(3,2,0)+(0,-1.2mm)$) {②};
  \node[font=\sffamily\small, text=bkGreen!50!black, anchor=east, inner sep=1pt] at ($(3,4,1.4)+(-1.2mm,0)$) {③};
  \node[font=\sffamily\footnotesize, anchor=north west, inner sep=1pt] at ($(3,4,0)+(0.8mm,-0.4mm)$) {Q};
\end{tikzpicture}
\end{minipage}\hfill
\begin{minipage}[c]{0.46\linewidth}
点P$(3,\ 4,\ 2)$ へは、原点Oから次の順に進みます。
\begin{enumerate}[label={\bkcircled[bkOrange]{\arabic*}}, leftmargin=2.2\zw, itemsep=1pt, topsep=2pt]
\item $x$ 軸にそって3（手前へ）
\item $y$ 軸の向きに4（右へ）
\item $z$ 軸の向きに2（上へ）
\end{enumerate}
3つの矢印は、たて・横・高さが3・4・2の箱（直方体）の辺です。Pは、Oと反対がわの箱の角にあります。
\end{minipage}

\par\vspace{4pt}
図の箱の点線は、見えない向こうがわの辺です。この本の図では、$x$ 軸を左下へななめにかいています。紙は平らなので、手前に出てくる向きを、ななめの線でかくのです。ななめにかいた $x$ 軸は、ほかの軸より短く見えるようにかいてあります（1目もりが、ほかの軸の約0.6倍）。目で長さをくらべないように気をつけましょう。

\bkivnstep{2}{反対に、点から3つの数を読む}
こんどは、空間の点Pが先にあって、その座標を読みとります。
\begin{enumerate}[label={\bkcircled[bkOrange]{\arabic*}}, leftmargin=2.4\zw, itemsep=1pt, topsep=2pt]
\item Pから床（$z=0$ の平面）へ、まっすぐ下に線をおろす。床にとどいた点をQとする。Qは、真上から照らしたときのPの「影」です。
\item Qは床の上の点なので、3-2と同じように、$x$ と $y$ を読む。図ではQ$(3,\ 4,\ 0)$。
\item PQの長さ（床からの高さ）が $z$。図では2。
\end{enumerate}
だから、Pの座標は $(3,\ 4,\ 2)$ です。進んで作っても、影から読んでも、同じ3つの数になりました。

このように、原点Oで直角に交わる3本の軸を決めると、
\bkivnline{\fbox{\textbf{空間の中の場所は、高さの $z$ を加えた3つの数 $(x,\ y,\ z)$ で表せる}}}
これを点の\textbf{座標}といい、P$(x,\ y,\ z)$ と書きます。\bkidx{くうかんのざひょう}{空間の座標}かっこの中の順番は、いつも $x$、$y$、$z$ です。$(3,\ 4,\ 2)$ と $(4,\ 3,\ 2)$ は別の点です。

\bkivnhead{\bkhatenaref の答え}
ドローンの場所を言うには、3つの数が必要です。たとえば部屋のすみを原点にして「手前へ\qty{2}{m}、右へ\qty{3}{m}、床から\qty{1.6}{m}の高さ」、つまり $(2,\ 3,\ 1.6)$ と伝えれば、友だちはドローンの場所を1つに決められます。電話の前に「原点はどのすみか、どっちを＋にするか」を、2人で決めておくことが大切です（3-2）。

\bkivnstep{3}{軸の上の点、座標平面の上の点}
3つの数のうち、いくつかが0になる点もあります。2本の軸をふくむ平面を、\textbf{座標平面}といいます。
\begin{bkfigure}
\begin{tikzpicture}[ivn3d=5.5mm]
  \foreach \sx/\nm/\cl/\eq in {0/xy/bkYellow/z=0, 1/yz/bkBlue/x=0, 2/zx/bkGreen/y=0} {
  \begin{scope}[xshift=\sx*46mm]
    \ifnum\sx=0 \fill[bkYellow!30] (0,0,0) -- (3.6,0,0) -- (3.6,4.4,0) -- (0,4.4,0) -- cycle; \fi
    \ifnum\sx=1 \fill[bkBlue!20] (0,0,0) -- (0,4.4,0) -- (0,4.4,3.4) -- (0,0,3.4) -- cycle; \fi
    \ifnum\sx=2 \fill[bkGreen!20] (0,0,0) -- (3.6,0,0) -- (3.6,0,3.4) -- (0,0,3.4) -- cycle; \fi
    \bkivnaxes{4.4}{5.2}{4.2}
    \node[font=\sffamily\small, anchor=north] at ($(0,2.2,0)+(0,-11mm)$) {{\bfseries $\nm$ 平面}（$\eq$）};
  \end{scope}}
\end{tikzpicture}
\bkfigcaption{$xy$ 平面（床）では $z=0$、$yz$ 平面（おくのかべ）では $x=0$、$zx$ 平面（左のかべ）では $y=0$。}
\end{bkfigure}
\begin{itemize}[leftmargin=1.2\zw, itemsep=2pt, topsep=2pt, label={\color{\bkTT}●}]
\item 床の上の点は、高さが0なので $z=0$。$xy$ 平面の上にあります。やってみようの「時計」は $y=0$ なので、$zx$ 平面（左のかべ）の上にあります。
\item $x$ 軸の上の点は、$y$ も $z$ も0で、$(3,\ 0,\ 0)$ のように書けます。
\item 3-2でマイナスの座標を使ったように、空間でもマイナスを使います。地下の駐車場なら、$z$ がマイナスになります。原点は $(0,\ 0,\ 0)$ です。
\end{itemize}

\bksubsection{2}{3本の軸の並べ方：右手系}

\leavevmode\label{bk:4-13:migite}3-2では、原点と＋の向きは自分で決めてよい、と学びました。空間でも同じです。ただし、空間には、もう1つだけ大事な約束があります。

\noindent\begin{minipage}[c]{0.5\linewidth}
右手の親指・人さし指・中指を、たがいに直角になるように開いてみましょう。
\begin{itemize}[leftmargin=1.2\zw, itemsep=1pt, topsep=2pt, label={\color{\bkTT}●}]
\item 親指 → $x$ 軸の＋の向き
\item 人さし指 → $y$ 軸の＋の向き
\item 中指 → $z$ 軸の＋の向き
\end{itemize}
この形になるように3本の軸を並べたものを、\textbf{右手系}（みぎてけい）といいます。\bkidx{みぎてけい}{右手系}この本の図は、すべて右手系です。
\end{minipage}\hfill
\begin{minipage}[c]{0.46\linewidth}
\centering
\begin{tikzpicture}[ivn3d=10mm]
  % 手のひら
  \fill[bkOrange!22] (0,0,0) ++(-4mm,-4.5mm) rectangle ++(8mm,7mm);
  \draw[bkOrange!70!black, line width=0.6pt, rounded corners=2mm] (0,0,0) ++(-4mm,-4.5mm) rectangle ++(8mm,7mm);
  \bkivnfinger{(0,0,0)}{(2.3,0,0)}{3.4mm}
  \bkivnfinger{(0,0,0)}{(0,2.4,0)}{3.2mm}
  \bkivnfinger{(0,0,0)}{(0,0,2.2)}{3.2mm}
  \draw[bkInk, line width=0.8pt, -{Stealth[length=4pt]}] (2.4,0,0) -- (3.4,0,0) node[anchor=north east, inner sep=1pt, font=\sffamily\small] {$x$};
  \draw[bkInk, line width=0.8pt, -{Stealth[length=4pt]}] (0,2.5,0) -- (0,3.3,0) node[anchor=west, inner sep=2pt, font=\sffamily\small] {$y$};
  \draw[bkInk, line width=0.8pt, -{Stealth[length=4pt]}] (0,0,2.3) -- (0,0,3.0) node[anchor=south, inner sep=2pt, font=\sffamily\small] {$z$};
  \node[font=\sffamily\footnotesize, anchor=east, inner sep=1pt] at ($(1.6,0,0)+(-2.4mm,1.6mm)$) {親指};
  \node[font=\sffamily\footnotesize, anchor=north, inner sep=1pt] at ($(0,1.5,0)+(0,-2.4mm)$) {人さし指};
  \node[font=\sffamily\footnotesize, anchor=west, inner sep=1pt] at ($(0,0,1.5)+(2.4mm,0)$) {中指};
\end{tikzpicture}
\end{minipage}
\par\vspace{6pt}
\bkivnhead{もう1つの見方：右ねじ}
ねじ（ペットボトルのふたや、くぎのようなねじ）は、上から見て右回り（時計回り）に回すと、おくへ進みます。このようなねじを\textbf{右ねじ}といいます。ふつうのねじは、ほとんどが右ねじです。

$x$ 軸を、$y$ 軸のほうへ90°回すことを考えます。このように回したときに、右ねじが進む向きが $z$ 軸の＋の向きです。$z$ 軸の上から見ると、$x$ から $y$ への回り方は左回り（反時計回り）です。下から見ると右回りなので、ねじは上へ進みます。

\begin{bkfigure}
\begin{tikzpicture}[ivn3d=8mm]
  % 右手系
  \bkivnaxes{3.6}{3.6}{3.2}
  \draw[bkRed, line width=1.1pt, -{Stealth[length=5pt]}] (1.5,0,0) .. controls (1.5,1.0,0) and (1.0,1.5,0) .. (0,1.5,0);
  \node[font=\sffamily\footnotesize, text=bkRed!80!black, anchor=north west, inner sep=1pt] at ($(1.2,1.2,0)+(0.5mm,-0.5mm)$) {$x$ から $y$ へ回す};
  \draw[bkGreen!60!black, line width=2.2pt, -{Stealth[length=6pt]}] (0,0,0.4) -- (0,0,2.1);
  \node[font=\sffamily\footnotesize, text=bkGreen!50!black, anchor=east, inner sep=1.5pt, align=right] at ($(0,0,1.4)+(-1.5mm,0)$) {右ねじが\\進む向き};
  \node[font=\sffamily\small\bfseries, anchor=north] at ($(0,1.6,0)+(0,-12mm)$) {右手系（この本の約束）};
  % 鏡
  \draw[bkGray, line width=1.2pt] ($(0,4.6,0)+(0,-8mm)$) -- ($(0,4.6,0)+(0,26mm)$);
  \node[font=\sffamily\footnotesize, text=bkGray, anchor=south] at ($(0,4.6,0)+(0,26.5mm)$) {鏡};
  % 左手系（鏡にうつしたもの）
  \begin{scope}[xshift=70mm]
  \draw[bkGray, line width=0.6pt, -{Stealth[length=4pt]}] (0,0,0) -- (0,-1.6,0);
  \draw[bkGray, line width=0.6pt, -{Stealth[length=4pt]}] (0,0,0) -- (3.6,0,0) node[anchor=north east, inner sep=1pt, font=\sffamily\small, text=bkInk] {$x$};
  \draw[bkGray, line width=0.6pt, -{Stealth[length=4pt]}] (0,0,0) -- (0,0,3.2) node[anchor=south, inner sep=2pt, font=\sffamily\small, text=bkInk] {$z$};
  \node[font=\sffamily\small, text=bkInk, anchor=east, inner sep=2pt] at (0,-1.6,0) {$y$};
  \node[font=\sffamily\small\bfseries, anchor=north, align=center] at ($(0,0.3,0)+(0,-12mm)$) {左手系（使わない）};
  \end{scope}
\end{tikzpicture}
\bkfigcaption{鏡にうつすと、$y$ 軸が反対を向いた「左手系」になる。左手系は、回しても右手系にはもどらない（右手の手ぶくろを、回しても左手の手ぶくろにできないのと同じ）。}
\end{bkfigure}

\bkivnhead{なぜ、この約束が大事なの？}
軸の並べ方を鏡にうつしたように変えると、$y$ の＋とマイナスが入れかわります。すると、次の章（4-14）で学ぶ「回す向き」の計算で、答えの符号（＋か－か）が反対になってしまいます。物理や工学の本の式は、すべて右手系で書かれています。設計図を作るソフトやロボットの制御も、右手系が基本です。だから、自分で軸を決めるときは、＋の向きと原点は自由に選んでよいが、\textbf{並べ方は右手系にする}、と約束します（ゲームを作るソフトの中には、左手系を使うものもあります）。

\bksubsection{3}{空間の2点の距離}

\leavevmode\label{bk:4-13:kyori}平面では、2点の距離を三平方の定理で求めました（4-4）。空間では、どうすればよいでしょう。

\bkivnstep{1}{箱に入る、いちばん長い棒}
横\qty{40}{cm}、おく行き\qty{30}{cm}、高さ\qty{20}{cm}の段ボール箱があります。長さ\qty{52}{cm}の棒は、この箱に入るかな？

\noindent\begin{minipage}[c]{0.44\linewidth}
\centering
\begin{tikzpicture}[ivn3d=1.0mm]
  \fill[bkOrange!12] (30,0,0) -- (30,40,0) -- (0,40,0) -- (0,0,0) -- cycle;
  \bkivnbox{30}{40}{20}{bkOrange!80!black}
  \draw[bkBlue, line width=1.4pt] (30,0,0) -- (0,40,0);
  \draw[bkGreen!70!black, line width=1.4pt] (0,40,0) -- (0,40,20);
  \draw[bkRed, line width=1.6pt] (30,0,0) -- (0,40,20);
  \draw[bkInk, line width=0.4pt] (0,40,2.5) -- ($(0,40,2.5)!0.1!(30,0,2.5)$) -- ($(0,40,0)!0.1!(30,0,0)$);
  \node[font=\sffamily\footnotesize, anchor=north, inner sep=1.5pt] at (30,20,0) {40};
  \node[font=\sffamily\footnotesize, anchor=north west, inner sep=1.5pt] at (15,40,0) {30};
  \node[font=\sffamily\footnotesize, anchor=west, inner sep=1.5pt] at (0,40,10) {20};
  \node[font=\sffamily\footnotesize, text=bkBlue!80!black, anchor=north east, inner sep=1pt] at (12,24,0) {50};
  \node[font=\sffamily\small, text=bkRed!80!black, anchor=south east, inner sep=1pt] at (15,20,10) {?};
\end{tikzpicture}
\end{minipage}\hfill
\begin{minipage}[c]{0.52\linewidth}
\textbf{① 床の対角線（青）}\par
底の長方形で、三平方の定理（4-4）。
\[ \sqrt{40^2+30^2}=\sqrt{2500}=50 \]
底にななめにおいても、\qty{50}{cm}までしか入りません。\par\vspace{3pt}
\textbf{② 箱の対角線（赤）}\par
青い線と、高さの線（緑）は直角です。もう一度、三平方の定理を使えます。
\[ \sqrt{50^2+20^2}=\sqrt{2900}\approx53.9 \]
\end{minipage}

\par\vspace{4pt}
$40^2+30^2=1600+900=2500$、$50^2+20^2=2500+400=2900$ です。箱の底のすみから、反対がわの天井のすみへ、ななめに入れると、\qty{53.9}{cm}まで入ります。だから\qty{52}{cm}の棒は、ななめにすれば入ります。三平方の定理を2回使うと、空間のななめの長さがわかるのです（4-4の直方体の対角線と同じ）。

\bkivnstep{2}{原点からの距離}
同じことを、座標でやってみます。原点Oと点P$(2,\ 3,\ 6)$ の距離を求めます。Pの真下の床の点をQ$(2,\ 3,\ 0)$ とします。

\noindent\begin{minipage}[c]{0.4\linewidth}
\centering
\begin{tikzpicture}[ivn3d=6mm]
  \bkivnaxes{3.6}{4.6}{7}
  \bkivnbox{2}{3}{6}{bkGray!70}
  \draw[bkBlue, line width=1.4pt] (0,0,0) -- (2,3,0);
  \draw[bkGreen!70!black, line width=1.4pt] (2,3,0) -- (2,3,6);
  \draw[bkRed, line width=1.6pt] (0,0,0) -- (2,3,6);
  \draw[bkInk, line width=0.4pt] (2,3,0.6) -- (1.76,2.64,0.6) -- (1.76,2.64,0);
  \fill[bkInk] (2,3,6) circle (0.8mm);
  \node[font=\sffamily\small, anchor=south west, inner sep=1.5pt] at (2,3,6) {P};
  \node[font=\sffamily\small, anchor=north west, inner sep=1.5pt] at (2,3,0) {Q};
  \node[font=\sffamily\footnotesize, anchor=east, inner sep=1.5pt] at (0,0,0) {O};
\end{tikzpicture}
\end{minipage}\hfill
\begin{minipage}[c]{0.56\linewidth}
\textbf{① 床の上で：OQ}\par
Qは床の上の点で、Oから手前へ2、右へ3。床の上の三平方の定理で、
\[ \mathrm{OQ}^2=2^2+3^2=4+9=13 \]
\textbf{② 立てた三角形で：OP}\par
PQは真上を向いているので、OQとPQは直角。三角形OQPで、もう一度三平方の定理。
\[ \mathrm{OP}^2=\mathrm{OQ}^2+\mathrm{PQ}^2=13+6^2=49 \]
$\mathrm{OP}=\sqrt{49}=7$
\end{minipage}

\par\vspace{4pt}
①の $\mathrm{OQ}^2=2^2+3^2$ を、②の式にそのまま入れると、$\mathrm{OP}^2=2^2+3^2+6^2$ です。途中の13を計算しなくても、3つの座標を2乗してたせばよいのです。だから、点P$(x,\ y,\ z)$ と原点の距離は
\bkivnline{\fbox{$\mathrm{OP}=\sqrt{x^2+y^2+z^2}$}}
です。平面の式 $\sqrt{x^2+y^2}$ に、高さの分の $z^2$ が1つふえただけです。

\bkivnstep{3}{2点のあいだの距離}
2点A$(x_1,\ y_1,\ z_1)$ とB$(x_2,\ y_2,\ z_2)$ のあいだでは、AからBまでの箱の、たて・横・高さが
\[ x_2-x_1,\qquad y_2-y_1,\qquad z_2-z_1 \]
になります（4-4の平面の2点の距離と同じで、座標の「あと－まえ」）。だから、
\bkivnline{\fbox{$\mathrm{AB}=\sqrt{(x_2-x_1)^2+(y_2-y_1)^2+(z_2-z_1)^2}$}}

\textbf{やってみようの部屋で}\quad 机の角A$(1.2,\ 0.8,\ 0.7)$ から電灯B$(2.0,\ 1.8,\ 2.5)$ までの距離は、
座標の差は、$x$ が $2.0-1.2=0.8$、$y$ が $1.8-0.8=1.0$、$z$ が $2.5-0.7=1.8$ なので、
\[ \mathrm{AB}=\sqrt{0.8^2+1.0^2+1.8^2}=\sqrt{0.64+1+3.24}=\sqrt{4.88}\approx2.21 \]
で、約\qty{2.21}{m}です。2乗するので、ひく順番を逆にして $1.2-2.0=-0.8$ としても、$(-0.8)^2=0.64$ で答えは同じです。

\bksubsection{4}{3次元のベクトル}

\leavevmode\label{bk:4-13:vector}ベクトルは「向きと大きさをもつ量」でした（4-10）。空を飛ぶドローンの動き・風・力は、上下の向きももっています。空間のベクトルも、平面と同じく成分で書けます。

\noindent\begin{minipage}[c]{0.44\linewidth}
\centering
\begin{tikzpicture}[ivn3d=3.2mm]
  \bkivnaxes{5.6}{6.4}{14}
  \bkivnbox{3}{4}{12}{bkGray!60}
  \draw[bkRed, line width=1.4pt] (0,0,0) -- (3,0,0);
  \draw[bkBlue, line width=1.4pt] (3,0,0) -- (3,4,0);
  \draw[bkGreen!70!black, line width=1.4pt] (3,4,0) -- (3,4,12);
  \draw[bkInk, line width=1.8pt, -{Stealth[length=6pt]}] (0,0,0) -- (3,4,12);
  \node[font=\sffamily\footnotesize, text=bkRed!80!black, anchor=east, inner sep=1.5pt] at ($(1.5,0,0)+(-1.2mm,0.6mm)$) {東へ3};
  \node[font=\sffamily\footnotesize, text=bkBlue!80!black, anchor=north, inner sep=1.5pt] at (3,2,0) {北へ4};
  \node[font=\sffamily\footnotesize, text=bkGreen!50!black, anchor=west, inner sep=1.5pt] at (3,4,7) {上へ12};
  \node[font=\sffamily\small, anchor=east, inner sep=2pt] at ($(1.5,2,6)+(-1mm,0)$) {$\vec a$};
\end{tikzpicture}
\end{minipage}\hfill
\begin{minipage}[c]{0.52\linewidth}
校庭から飛び立ったドローンが、東へ\qty{3}{m}、北へ\qty{4}{m}、上へ\qty{12}{m}動きました。東を $x$、北を $y$、上を $z$ の＋の向きにします（この並べ方は右手系です）。

この動き（変位）のベクトル $\vec a$ は、3つの成分で
\[ \vec a=(3,\ 4,\ 12) \]
と書けます。\textbf{成分}は、それぞれの軸の向きに、どれだけ進むかを表します（4-11）。
\end{minipage}

\par\vspace{4pt}
\begin{itemize}[leftmargin=1.2\zw, itemsep=3pt, topsep=2pt, label={\color{\bkTT}●}]
\item \textbf{大きさ}\quad 4-11では、成分が3つ以上のときも、大きさを「成分の2乗の和の $\sqrt{\ }$」と決めました。その決め方が図の長さと合うことは、③の箱の対角線で確かめたとおりです。床の対角線（三平方の定理）と高さで、もう一度三平方の定理を使うと、ちょうど $\sqrt{a_x^2+a_y^2+a_z^2}$ になりました。\\
$|\vec a|=\sqrt{3^2+4^2+12^2}=\sqrt{9+16+144}=\sqrt{169}=13$。ドローンは、まっすぐ\qty{13}{m}動いたことになります。
\item \textbf{2点を結ぶベクトル}\quad 点A$(x_1,\ y_1,\ z_1)$ からB$(x_2,\ y_2,\ z_2)$ へのベクトルは、成分が「あと－まえ」です。\\
$\overrightarrow{\mathrm{AB}}=(x_2-x_1,\ y_2-y_1,\ z_2-z_1)$
\item \textbf{和・差・実数倍は、成分ごと}\quad 平面と同じです（4-11）。$x$ は $x$ どうし、$y$ は $y$ どうし、$z$ は $z$ どうしで計算します。たとえば $\vec a=(1,\ 2,\ 2)$、$\vec b=(3,\ 0,\ -4)$ なら\\
$\vec a+\vec b=(1+3,\ 2+0,\ 2+(-4))=(4,\ 2,\ -2)$、\quad $3\vec a=(3,\ 6,\ 6)$
\end{itemize}

\bkivnhead{動く点の位置：速度をかけてたす}
ドローンが、高さ\qty{10}{m}の点 $\vec p_0=(0,\ 0,\ 10)$ から、速度 $\vec v=(2,\ 1,\ 2)$（m/s）で、まっすぐ同じ速さで飛びます。1秒ごとに、東へ2、北へ1、上へ2ずつ進みます。速さ（速度の大きさ）は $|\vec v|=\sqrt{4+1+4}=3$ で、\qty{3}{m/s}です。$t$ 秒後の位置は、4-11の等速直線運動の式 $\vec p=\vec p_0+t\vec v$ と同じで、
\[ \vec p=(0,\ 0,\ 10)+t\,(2,\ 1,\ 2)=(2t,\ t,\ 10+2t) \]
5秒後は $(10,\ 5,\ 20)$ です。成分が3つにふえても、ベクトルの計算のきまりは、まったく同じです（4-11の「数の組は、成分が3つでも10個でもベクトル」）。

\bksubsection{5}{3次元の内積となす角}

\leavevmode\label{bk:4-13:naiseki}2つのベクトル $\vec a$、$\vec b$ の内積は、なす角を $\theta$ として $\vec a\cdot\vec b=|\vec a||\vec b|\cos\theta$ でした（4-12）。空間のベクトルでも、なす角を考えられるでしょうか。

\bkivnstep{1}{2本の矢印は、いつも1枚の平面にのる}
えんぴつを2本、おしりをくっつけて、ちがう向きに持ってみましょう。下じきを当てると、2本のえんぴつに同時にふれる向きが、必ずあります。空間の2つのベクトルも、根もとをそろえると、その2本をふくむ平面が1枚決まります。その平面の上で見れば、2本の矢印のあいだの角 $\theta$ は、平面のときと同じようにはかれます。だから、なす角も内積も、空間でそのまま使えます。

\bkivnstep{2}{成分の式も、$z$ の分がふえるだけ}
平面では、余弦定理（4-7）から $\vec a\cdot\vec b=a_xb_x+a_yb_y$ を作りました（4-12）。空間でも同じように作ります。$\vec a$、$\vec b$ と、その先を結ぶ $\vec b-\vec a$ で三角形ができます。この三角形も1枚の平面の上にあるので、余弦定理が使えます。
\[ |\vec b-\vec a|^2=|\vec a|^2+|\vec b|^2-2|\vec a||\vec b|\cos\theta\qquad\cdots\text{（ア）} \]
左がわを、成分で計算します。$(b_x-a_x)^2=b_x^2-2a_xb_x+a_x^2$（3-8の展開）なので、
\begin{align*}
|\vec b-\vec a|^2&=(b_x-a_x)^2+(b_y-a_y)^2+(b_z-a_z)^2\\
&=(a_x^2+a_y^2+a_z^2)+(b_x^2+b_y^2+b_z^2)-2(a_xb_x+a_yb_y+a_zb_z)\\
&=|\vec a|^2+|\vec b|^2-2(a_xb_x+a_yb_y+a_zb_z)
\end{align*}
（ア）の右がわとくらべると、$|\vec a|^2+|\vec b|^2$ は同じなので、のこりの部分も等しいはずです。
\bkivnline{\fbox{$\vec a\cdot\vec b=|\vec a||\vec b|\cos\theta=a_xb_x+a_yb_y+a_zb_z$}}
平面の式に、$z$ どうしをかけた $a_zb_z$ が1つふえただけです。

\bkivnstep{3}{図と計算で、たしかめる}
\noindent\begin{minipage}[c]{0.4\linewidth}
\centering
\begin{tikzpicture}[ivn3d=22mm]
  \bkivnbox{1}{1}{1}{bkGray!70}
  \draw[bkBlue, line width=1.5pt, -{Stealth[length=5pt]}] (0,0,0) -- (1,1,0);
  \draw[bkOrange, line width=1.5pt, -{Stealth[length=5pt]}] (0,0,0) -- (0,1,1);
  \draw[bkGreen!70!black, line width=1pt, dashed] (1,1,0) -- (0,1,1);
  \fill[bkInk] (0,0,0) circle (0.7mm);
  % ラベルは、立方体の辺・矢印から 1.5 mm 以上はなれた空いている所に置く（第14回 B7）
  \node[font=\sffamily\footnotesize, anchor=east, inner sep=1pt] at ($(0,0,0)+(-1.2mm,1.2mm)$) {O};
  \node[font=\sffamily\footnotesize, text=bkBlue!80!black, inner sep=0pt] at ($(0,0,0)+(-0.3mm,-4.2mm)$) {$\vec a$};
  \node[font=\sffamily\footnotesize, text=bkOrange!80!black, anchor=south east, inner sep=0pt] at ($(0,0.35,0.35)+(-1.5mm,1.5mm)$) {$\vec b$};
  \node[font=\sffamily\footnotesize, anchor=north, inner sep=0pt] at ($(1,1,0)+(0,-1.8mm)$) {(1, 1, 0)};
  \node[font=\sffamily\footnotesize, anchor=south, inner sep=0pt] at ($(0,1,1)+(0,1.8mm)$) {(0, 1, 1)};
\end{tikzpicture}
\end{minipage}\hfill
\begin{minipage}[c]{0.56\linewidth}
1辺が1の立方体で、$\vec a=(1,\ 1,\ 0)$（床の面の対角線）と $\vec b=(0,\ 1,\ 1)$（おくのかべの面の対角線）のなす角を求めます。\par\vspace{3pt}
\textbf{図で}\quad 3つの辺（青・オレンジ・緑の点線）は、どれも立方体の面の対角線で、長さは $\sqrt{1^2+1^2}=\sqrt2$。3つの辺が等しいので正三角形で、なす角は60°です。\par\vspace{3pt}
\textbf{計算で}\quad $\vec a\cdot\vec b=1\times0+1\times1+0\times1=1$\\
$|\vec a||\vec b|=\sqrt2\times\sqrt2=2$ なので\\
$\cos\theta=\dfrac{1}{2}$、$\theta=60°$
\end{minipage}

\par\vspace{4pt}
図で考えた60°と、成分の式で計算した60°が、ぴったり合いました。

\bkivnhead{内積の使い道（平面と同じ）}
\begin{itemize}[leftmargin=1.2\zw, itemsep=2pt, topsep=2pt, label={\color{\bkTT}●}]
\item \textbf{なす角}\quad $\cos\theta=\dfrac{\vec a\cdot\vec b}{|\vec a||\vec b|}$ を計算し、電卓の $\cos^{-1}$ で角度にする（4-6）。電卓は DEG（度）にしておく。
\item \textbf{直交}\quad $\vec a\cdot\vec b=0$ なら、2つのベクトルは直角（$\cos90°=0$）。
\item \textbf{仕事}\quad 力 $\vec F$ で、物が $\vec s$ だけ動いたときの仕事は（4-12）\\ $W=\vec F\cdot\vec s=F_xs_x+F_ys_y+F_zs_z$\\ 上向きの力は、物を持ち上げる分だけ仕事をします。
\end{itemize}

\bksubsection{6}{基本ベクトルと、向きを表す単位ベクトル}

\bkivnstep{1}{3本の軸の向きの、長さ1の矢印}\label{bk:4-13:kihon}
$x$ 軸・$y$ 軸・$z$ 軸の＋の向きを向いた、長さ1のベクトルを
\[ \vec i=(1,\ 0,\ 0),\qquad \vec j=(0,\ 1,\ 0),\qquad \vec k=(0,\ 0,\ 1) \]
と書き、\textbf{基本ベクトル}といいます。\bkidx{きほんべくとる}{基本ベクトル}どれも大きさが1の単位ベクトル（4-11）で、たがいに直交しています。

4-11では、$x$ の向きの単位ベクトルを $\vec e_x=(1,\ 0)$、$y$ の向きを $\vec e_y=(0,\ 1)$ と書きました。$\vec i$、$\vec j$ は、これと同じものの別の書き方です。空間では $z$ の向きの $\vec e_z$ が加わり、
\[ \vec e_x=\vec i,\qquad \vec e_y=\vec j,\qquad \vec e_z=\vec k \]
です。どちらの書き方も使われるので、両方読めるようにしておきましょう（平面では $\vec i$ と $\vec j$ の2本だけを使います）。

\noindent\begin{minipage}[c]{0.44\linewidth}
\centering
\begin{tikzpicture}[ivn3d=7mm]
  \bkivnaxes{4.6}{5.4}{4.4}
  \bkivnbox{2}{3}{3}{bkGray!60}
  \draw[bkRed, line width=1.6pt, -{Stealth[length=4pt]}] (0,0,0) -- (1,0,0);
  \draw[bkRed, line width=1.6pt, -{Stealth[length=4pt]}] (1,0,0) -- (2,0,0);
  \draw[bkBlue, line width=1.6pt, -{Stealth[length=4pt]}] (2,0,0) -- (2,1,0);
  \draw[bkBlue, line width=1.6pt, -{Stealth[length=4pt]}] (2,1,0) -- (2,2,0);
  \draw[bkBlue, line width=1.6pt, -{Stealth[length=4pt]}] (2,2,0) -- (2,3,0);
  \draw[bkGreen!70!black, line width=1.6pt, -{Stealth[length=4pt]}] (2,3,0) -- (2,3,1);
  \draw[bkGreen!70!black, line width=1.6pt, -{Stealth[length=4pt]}] (2,3,1) -- (2,3,2);
  \draw[bkGreen!70!black, line width=1.6pt, -{Stealth[length=4pt]}] (2,3,2) -- (2,3,3);
  \draw[bkInk, line width=1.2pt, -{Stealth[length=5pt]}] (0,0,0) -- (2,3,3);
  \node[font=\sffamily\footnotesize, text=bkRed!80!black, anchor=north east, inner sep=1pt] at (1.2,0,0) {$\vec i$ が2本};
  \node[font=\sffamily\footnotesize, text=bkBlue!80!black, anchor=north, inner sep=2pt] at (2,1.5,0) {$\vec j$ が3本};
  \node[font=\sffamily\footnotesize, text=bkGreen!50!black, anchor=west, inner sep=2pt] at (2,3,1.5) {$\vec k$ が3本};
\end{tikzpicture}
\end{minipage}\hfill
\begin{minipage}[c]{0.52\linewidth}
ベクトル $\vec a=(2,\ 3,\ 3)$ は、$\vec i$ を2本、$\vec j$ を3本、$\vec k$ を3本つないだ矢印と同じです。だから
\[ \vec a=2\vec i+3\vec j+3\vec k \]
と書けます。成分は「基本ベクトルが何本分か」を表す数です。いっぱんに、
\[ \vec a=(a_x,\ a_y,\ a_z)=a_x\vec i+a_y\vec j+a_z\vec k \]
「東へ2歩、北へ3歩、上へ3段」と言うのと同じです。
\end{minipage}

\par\vspace{4pt}
内積を使うと、成分を取り出せます。$\vec a\cdot\vec i=a_x\times1+a_y\times0+a_z\times0=a_x$ です。「$\vec a$ が $x$ 軸の向きにどれだけ進むか」が、$\vec i$ との内積で求まります（4-12の正射影）。

\bkivnstep{2}{向きだけを取り出す：単位ベクトルと方向余弦}
ベクトル $\vec a$ を、その大きさ $|\vec a|$ でわると、向きは同じで、長さ1のベクトルになります（4-11の単位ベクトル）。$\vec a=(2,\ 3,\ 6)$ なら、$|\vec a|=\sqrt{4+9+36}=7$ なので、
\[ \vec e=\frac{\vec a}{|\vec a|}=\Bigl(\frac{2}{7},\ \frac{3}{7},\ \frac{6}{7}\Bigr) \]
この3つの数には、意味があります。$\vec a$ と $x$ 軸のなす角を $\theta_x$ とすると、⑤の式で
\[ \cos\theta_x=\frac{\vec a\cdot\vec i}{|\vec a||\vec i|}=\frac{a_x}{|\vec a|}=\frac{2}{7} \]
$y$ 軸・$z$ 軸とのなす角 $\theta_y$、$\theta_z$ も同じで、$\cos\theta_y=\dfrac{3}{7}$、$\cos\theta_z=\dfrac{6}{7}$。つまり、
\bkivnline{\fbox{$\vec e=\dfrac{\vec a}{|\vec a|}=(\cos\theta_x,\ \cos\theta_y,\ \cos\theta_z)$}}
単位ベクトルの3つの成分は、3本の軸とのなす角の cos です。これを\textbf{方向余弦}（ほうこうよげん）といいます。\bkidx{ほうこうよげん}{方向余弦}電卓の $\cos^{-1}$ で角度にすると、
\begin{bkcenter}
\begin{bktable}{colspec={Q[l,m,wd=22mm] *{3}{Q[c,m,wd=28mm]}}, rowsep=1.5pt, row{1}={font=\sffamily\small}}
 & $x$ 軸と & $y$ 軸と & $z$ 軸と \\
cos & $\dfrac{2}{7}\approx0.286$ & $\dfrac{3}{7}\approx0.429$ & $\dfrac{6}{7}\approx0.857$ \\
角 & $\theta_x\approx73.4°$ & $\theta_y\approx64.6°$ & $\theta_z\approx31.0°$ \\
\end{bktable}
\bktabcaption{$z$ 軸（真上）とのなす角がいちばん小さい。$\vec a$ は、上向きに近い矢印。}
\end{bkcenter}
$\vec e$ の長さは1なので、成分を2乗してたすと1になります。\[ \Bigl(\frac{2}{7}\Bigr)^2+\Bigl(\frac{3}{7}\Bigr)^2+\Bigl(\frac{6}{7}\Bigr)^2=\frac{4+9+36}{49}=\frac{49}{49}=1 \] つまり
\[ \cos^2\theta_x+\cos^2\theta_y+\cos^2\theta_z=1 \]
平面なら、$x$ 軸から角 $\theta$ の向きの単位ベクトルは $(\cos\theta,\ \sin\theta)$ でした（4-8）。$y$ 軸とのなす角は $90°-\theta$ なので、$\sin\theta=\cos(90°-\theta)$（4-6）を使うと、平面の単位ベクトルも「2本の軸とのなす角の cos」の組です。$\sin^2\theta+\cos^2\theta=1$ の3次元版が、上の式です。

\bksubsection{7}{空間の直線と平面}

\bkivnstep{1}{交わらないのに平行でもない：ねじれの位置}\label{bk:4-13:heimen}
平面の上の2本の直線は、「交わる」か「平行」のどちらかでした。空間では、もう1つの場合があります。

\noindent\begin{minipage}[c]{0.46\linewidth}
\centering
\begin{tikzpicture}[ivn3d=7mm]
  % 下の道路（x の向き）と、高架の道路（y の向き、高さ2）
  \fill[bkGray!22] (-1,2,0) -- (6,2,0) -- (6,3,0) -- (-1,3,0) -- cycle;
  \draw[bkInk!60, line width=0.5pt] (-1,2,0) -- (6,2,0) (-1,3,0) -- (6,3,0);
  \draw[bkBlue, line width=1.4pt] (-1,2.5,0) -- (6,2.5,0);
  \fill[bkGray!45] (2.5,0.6,0) -- (2.5,0.6,2) -- (2.1,0.6,2) -- (2.1,0.6,0) -- cycle;
  \fill[bkGray!45] (2.5,4.4,0) -- (2.5,4.4,2) -- (2.1,4.4,2) -- (2.1,4.4,0) -- cycle;
  \fill[bkGray!30] (3,-0.6,2) -- (3,5.6,2) -- (2,5.6,2) -- (2,-0.6,2) -- cycle;
  \draw[bkInk!60, line width=0.5pt] (3,-0.6,2) -- (3,5.6,2) (2,-0.6,2) -- (2,5.6,2);
  \draw[bkOrange, line width=1.4pt] (2.5,-0.6,2) -- (2.5,5.6,2);
  \node[font=\sffamily\footnotesize, text=bkBlue!80!black, anchor=north west, inner sep=1.5pt] at (6,3,0) {下の道路};
  \node[font=\sffamily\footnotesize, text=bkOrange!80!black, anchor=south, inner sep=2pt] at (2,4.6,2) {高架の道路};
\end{tikzpicture}
\end{minipage}\hfill
\begin{minipage}[c]{0.5\linewidth}
立体交差では、下の道路（青）と高架の道路（オレンジ）は、上から見ると十字に交わっています。でも、高さがちがうので、ぶつかりません。向きもちがうので、平行でもありません。

このように、\textbf{交わらず、平行でもない}2本の直線の関係を、\textbf{ねじれの位置}といいます。ねじれの位置は、平面の上では起こりません。空間ならではの関係です。
\end{minipage}

\bkivnstep{2}{平面の向きは、1本の矢印で決まる：法線ベクトル}
画びょうを板にまっすぐさすと、針は板に垂直です。板をかたむけると、針の向きも変わります。反対に、針の向きを決めると、板の向きが1つに決まります。

平面に垂直なベクトルを、その平面の\textbf{法線ベクトル}（ほうせんべくとる）といいます。\bkidx{ほうせんべくとる}{法線ベクトル}平面の上のどんな線とも直角に交わる矢印です。床（$xy$ 平面）の法線ベクトルは、真上を向いた $\vec k=(0,\ 0,\ 1)$ です。

\noindent\begin{minipage}[c]{0.46\linewidth}
\centering
\begin{tikzpicture}[ivn3d=6.2mm]
  \bkivnaxes{4.2}{6.8}{4}
  \filldraw[fill=bkTheme!25, draw=bkDeep, line width=0.8pt, opacity=0.9] (3,0,0) -- (0,6,0) -- (0,0,3) -- cycle;
  \fill[bkInk] (1,2,1) circle (0.8mm);
  \node[font=\sffamily\footnotesize, anchor=north, inner sep=1.5pt] at ($(1,2,1)+(0,-0.6mm)$) {P$_0$};
  \draw[bkRed, line width=1.6pt, -{Stealth[length=5pt]}] (1,2,1) -- (2.2,2.6,2.2);
  \node[font=\sffamily\small, text=bkRed!80!black, anchor=west, inner sep=2pt] at (2.2,2.6,2.2) {$\vec n=(2,\ 1,\ 2)$};
  \fill[bkInk] (2,0,1) circle (0.7mm);
  \draw[bkBlue, line width=1.1pt, -{Stealth[length=4pt]}] (1,2,1) -- (2,0,1);
  \node[font=\sffamily\footnotesize, anchor=north east, inner sep=1pt] at (2,0,1) {P};
  \node[font=\sffamily\footnotesize, anchor=north, inner sep=1.5pt] at (3,0,0) {3};
  \node[font=\sffamily\footnotesize, anchor=north, inner sep=1.5pt] at (0,6,0) {6};
  \node[font=\sffamily\footnotesize, anchor=east, inner sep=1.5pt] at (0,0,3) {3};
\end{tikzpicture}
\end{minipage}\hfill
\begin{minipage}[c]{0.5\linewidth}
点P$_0$ をふくみ、$\vec n$ に垂直な平面を考えます。平面の上のどの点Pでも、平面の上の矢印 $\overrightarrow{\mathrm{P_0P}}$（青）は $\vec n$ と直角です。直角なら内積は0（⑤）なので、
\[ \vec n\cdot\overrightarrow{\mathrm{P_0P}}=0 \]
これが、「Pがこの平面の上にある」という条件です。
\end{minipage}

\par\vspace{4pt}
図の平面で計算してみます。P$_0=(1,\ 2,\ 1)$、$\vec n=(2,\ 1,\ 2)$ とし、平面の上の点を P$(x,\ y,\ z)$ とします。$\overrightarrow{\mathrm{P_0P}}=(x-1,\ y-2,\ z-1)$ なので、
\begin{align*}
2(x-1)+1\times(y-2)+2(z-1)&=0 && \text{（内積が0）}\\
2x-2+y-2+2z-2&=0 && \text{（かっこをはずす）}\\
2x+y+2z&=6 && \text{（数を右へ移す）}
\end{align*}
これが、この平面の式です。\textbf{$x$、$y$、$z$ の前の数 $2,\ 1,\ 2$ は、法線ベクトルの成分そのもの}です。たしかめます。図の点P$(2,\ 0,\ 1)$ は、$2\times2+0+2\times1=6$ で式をみたすので、平面の上にあります。$x$ 軸との交わる点は、$y=z=0$ とおいて $2x=6$ から $x=3$ です（図の3）。いっぱんに、
\bkivnline{\fbox{$\vec n=(a,\ b,\ c)$ に垂直な平面の式は\enspace $ax+by+cz=d$}}
$d$ は、平面の上の1点の座標を入れて決めます。
平面の向きを法線ベクトルで表す考え方は、8-8で、平面を通りぬける水や電気の量を計算するときに使います。

\begin{motto}[4]
\textbf{原点から平面までの距離}\par
②の平面 $2x+y+2z=6$ と原点Oの距離（Oから平面に下ろした垂線の長さ）を求めます。法線ベクトル $\vec n=(2,\ 1,\ 2)$ の大きさは $\sqrt{4+1+4}=3$ なので、同じ向きの単位ベクトルは $\vec e=\Bigl(\dfrac{2}{3},\ \dfrac{1}{3},\ \dfrac{2}{3}\Bigr)$ です（⑥）。

平面の上のどの点Pでも、$\overrightarrow{\mathrm{OP}}$ を $\vec e$ の向きに正射影した長さ（4-12）は、Oから平面までの距離と同じです。正射影の長さは $\overrightarrow{\mathrm{OP}}\cdot\vec e$ なので、
\[ \overrightarrow{\mathrm{OP}}\cdot\vec e=\frac{2x+y+2z}{3}=\frac{6}{3}=2 \]
Pが平面のどこにあっても、$2x+y+2z$ はいつも6なので、答えは同じ2です。たしかめ：Oから $\vec e$ の向きに2進んだ点は $\Bigl(\dfrac{4}{3},\ \dfrac{2}{3},\ \dfrac{4}{3}\Bigr)$ で、$\dfrac{8}{3}+\dfrac{2}{3}+\dfrac{8}{3}=\dfrac{18}{3}=6$。ちゃんと平面の上にあります。いっぱんに、原点と平面 $ax+by+cz=d$ の距離は $\dfrac{|d|}{\sqrt{a^2+b^2+c^2}}$ です。
\end{motto}

\bksubsection{8}{空間の形を紙にかく：投影図}

\leavevmode\label{bk:4-13:toei}ここまでの図は、空間をななめから見た図（見取図）でした。ものを作るための設計図では、長さを正しく読めるように、真正面・真上・真横から見た図をかきます。これを\textbf{投影図}（とうえいず）といいます。

\begin{bkfigure}
\begin{tikzpicture}[x=1mm, y=1mm]
  % 見取図：L字のブロック。x(手前)0〜2、y(右)0〜4、z(上)：y=0〜1.5 の所は 0〜3、y=1.5〜4 の所は 0〜1（1目もり 7mm）
  \begin{scope}[shift={(4mm,6mm)}, ivn3d=7mm]
    \bkivnaxes{3.4}{5.2}{4.2}
    \filldraw[fill=bkBlue!14, draw=bkInk, line width=0.7pt] (2,0,0) -- (2,4,0) -- (2,4,1) -- (2,1.5,1) -- (2,1.5,3) -- (2,0,3) -- cycle;
    \filldraw[fill=bkOrange!18, draw=bkInk, line width=0.7pt] (2,0,3) -- (2,1.5,3) -- (0,1.5,3) -- (0,0,3) -- cycle;
    \filldraw[fill=bkGreen!18, draw=bkInk, line width=0.7pt] (2,1.5,3) -- (2,1.5,1) -- (0,1.5,1) -- (0,1.5,3) -- cycle;
    \filldraw[fill=bkOrange!18, draw=bkInk, line width=0.7pt] (2,1.5,1) -- (2,4,1) -- (0,4,1) -- (0,1.5,1) -- cycle;
    \filldraw[fill=bkGreen!18, draw=bkInk, line width=0.7pt] (2,4,0) -- (2,4,1) -- (0,4,1) -- (0,4,0) -- cycle;
    \node[font=\sffamily\small\bfseries, anchor=north] at ($(0,2,0)+(0,-9mm)$) {見取図};
  \end{scope}
  % 正面図（手前 +x から見る：横に y、たてに z）
  \begin{scope}[shift={(62,4)}, x=5mm, y=5mm]
    \filldraw[fill=bkBlue!14, draw=bkInk, line width=0.8pt] (0,0) -- (4,0) -- (4,1) -- (1.5,1) -- (1.5,3) -- (0,3) -- cycle;
    \node[font=\sffamily\footnotesize, anchor=north] at (2,-0.3) {正面図（手前から）};
  \end{scope}
  % 平面図（真上 +z から見る：横に y、たてに −x）
  \begin{scope}[shift={(62,30)}, x=5mm, y=5mm]
    \filldraw[fill=bkOrange!18, draw=bkInk, line width=0.8pt] (0,0) rectangle (4,2);
    \draw[bkInk, line width=0.8pt] (1.5,0) -- (1.5,2);
    \node[font=\sffamily\footnotesize, anchor=south] at (2,2.3) {平面図（真上から）};
  \end{scope}
  % 側面図（右 +y から見る：横に −x、たてに z）
  \begin{scope}[shift={(96,4)}, x=5mm, y=5mm]
    \filldraw[fill=bkGreen!18, draw=bkInk, line width=0.8pt] (0,0) rectangle (2,3);
    \draw[bkInk, line width=0.8pt] (0,1) -- (2,1);
    \node[font=\sffamily\footnotesize, anchor=north] at (1,-0.3) {側面図（右から）};
  \end{scope}
  \draw[bkGray, line width=0.4pt, dashed] (62,19.5) -- (62,30) (82,9) -- (82,30);
  \draw[bkGray, line width=0.4pt, dashed] (82,9) -- (96,9) (69.5,19) -- (96,19);
\end{tikzpicture}
\bkfigcaption{L字の形のブロックの見取図と投影図。正面図の真上に平面図、右に側面図を並べる（日本の設計図のきまり、第三角法）。点線は、同じ所をつなぐ線。}
\end{bkfigure}

座標で考えると、投影図は「1つの座標を消した図」です。
\begin{itemize}[leftmargin=1.2\zw, itemsep=2pt, topsep=2pt, label={\color{\bkTT}●}]
\item \textbf{平面図}（真上から見る）：高さ $z$ が消えて、$x$ と $y$ が残る。①の「影」と同じです。
\item \textbf{正面図}（手前から見る）：手前への $x$ が消えて、$y$ と $z$ が残る。
\item \textbf{側面図}（右から見る）：右への $y$ が消えて、$x$ と $z$ が残る。
\end{itemize}
1つの図では、1つの座標がわかりません。正面図だけでは、ブロックの奥行きはわかりません。だから設計図では、2つか3つの図を組にして、3つの座標をすべて読みとれるようにします。たとえば点P$(2,\ 0,\ 3)$（ブロックの左上の手前の角）は、正面図では左はしの、高さ3の所にかかれます。平面図では、いちばん下の辺（手前）の左はしにかかれます。

\bksubsection{9}{専門書の書き方：太字のベクトル}

\leavevmode\label{bk:4-13:futoji}この本の第4巻では、ベクトルを矢印つきの文字 $\vec a$ で書いてきました。大学の教科書など、工学の専門書では、矢印をつけずに、\textbf{太い字}でベクトルを表すことが多いです。このシリーズ（第7巻・第8巻）と物理シリーズの本文では、この本と同じ矢印 $\vec a$ で書き、専門書の書き方は「専門書をのぞいてみよう」で見せます。専門書を読むときのために、太字の読み方を知っておきましょう。
\begin{bkcenter}
\begin{bktable}{colspec={Q[l,m,wd=36mm] Q[c,m,wd=36mm] Q[c,m,wd=40mm]}, rowsep=2pt, row{1}={font=\sffamily\small}}
 & この巻の書き方 & 専門書の書き方 \\
ベクトル & $\vec a$、$\vec F$、$\overrightarrow{\mathrm{OP}}$ & $\symbf a$、$\symbf F$、$\symbf r$ \\
大きさ & $|\vec a|$ & $|\symbf a|$ または \textbf{細い字の} $a$ \\
基本ベクトル & $\vec i,\ \vec j,\ \vec k$ & $\symbf i,\ \symbf j,\ \symbf k$ や $\symbf e_x,\ \symbf e_y,\ \symbf e_z$ \\
成分で & $\vec a=(a_x,\ a_y,\ a_z)$ & $\symbf a=(a_x,\ a_y,\ a_z)$ \\
\end{bktable}
\bktabcaption{太字の $\symbf a$ はベクトル、細い字の $a$ はその大きさ（ふつうの数）。}
\end{bkcenter}
とくに大事なのは、\textbf{太い字の $\symbf a$ と細い字の $a$ は別のもの}だということです。物理の本で「$\symbf F$ の大きさを $F$ とする」と書いてあれば、$F=|\symbf F|$ のことです。力の向きを考えるときは $\symbf F$、大きさだけを計算するときは $F$ を使います。手でノートに書くときは、太い字は書きにくいので、矢印をつけて $\vec a$ と書いてかまいません。

\begin{senmonsho}[空間ベクトルの書き方]
\begin{bkhonbun}
位置ベクトル $\symbf r$ の成分表示と大きさは
\[ \symbf r=x\symbf i+y\symbf j+z\symbf k,\qquad r=|\symbf r|=\sqrt{x^2+y^2+z^2} \]
である。$\symbf a$ と $\symbf b$ のなす角を $\theta$ とすると
\[ \symbf a\cdot\symbf b=ab\cos\theta=a_xb_x+a_yb_y+a_zb_z \]
$\symbf a$ の向きの単位ベクトルは $\symbf a/a$ で、その成分を方向余弦という。
\end{bkhonbun}
\bkyaku{点の位置を表す矢印（位置ベクトル）を太字の $\symbf r$ で書く。その長さは細い字の $r$。$a$、$b$ は $\symbf a$、$\symbf b$ の大きさなので、$ab\cos\theta$ は $|\vec a||\vec b|\cos\theta$ のこと。$\symbf a/a$ は「$\symbf a$ をその大きさでわったもの」。}
\end{senmonsho}

\begin{motto}[3]
\textbf{4次元の空間ってあるの？}\par
友だちと待ち合わせをするとき、「駅ビルの3階の東口に」だけでは足りません。「午後3時に」という時刻もいります。場所の3つの数と、時刻の1つの数で、4つの数の組 $(x,\ y,\ z,\ t)$ になります。物理学では、これを「4次元の時空」とよびます。

数学では、数の組は成分がいくつでもベクトルと考えられました（4-11）。たとえば5人の体重の組も、5次元のベクトルです。長さも、成分を2乗してたして $\sqrt{\ }$ をとれば、何次元でも同じように計算できます。3-15の相関係数は、データの組を「たくさんの次元をもつベクトル」と見て、そのなす角の cos を計算していたのです（4-12）。目には見えない高い次元の空間も、数の組と計算のきまりで、3次元と同じようにあつかえます。
\end{motto}

\begin{column}{ハエが飛び立つと、数は3つ}
3-2のコラムで、天井をはうハエを見て座標を思いついたという、デカルトの言い伝えを読みました。天井の上なら2つの数で足りましたが、ハエが飛び立てば、床からの高さを加えた3つの数 $(x,\ y,\ z)$ がいります。空を飛ぶ飛行機も同じです。空港の管制では、地図の上の位置に「高度」を加えて、飛行機どうしが近づきすぎないように見守っています。上空を飛ぶ旅客機は、高さ\qty{10}{km}前後を飛び、同じ空の道でも、高さをずらして安全な間かくをとっています。
\end{column}

\begin{column}{日本とヨーロッパで、図の並べ方がちがう}
⑧の投影図では、正面図の上に平面図、右に側面図を並べました。これは「第三角法」というきまりで、日本の設計図（JIS）やアメリカでは、ふつうこの並べ方を使います。ところが、ヨーロッパでは、正面図の下に平面図、左に右側面図を並べる「第一角法」もよく使われます。同じ部品でも、図の並べ方を取りちがえると、上下や左右が反対の物ができてしまいます。そこで、設計図のすみには、どちらの並べ方でかいたかを示す小さなマークをかいておく約束になっています。
\end{column}

\begin{taisetsu}
\def\tsl#1{\makebox[5.8\zw][l]{\sffamily\bfseries\small\color{\bkTT}#1}}
\tsl{空間の座標}点は3つの数 $(x,\ y,\ z)$ で表す。軸は右手系に並べる。\\[3pt]
\tsl{距離}$\mathrm{AB}=\sqrt{(x_2-x_1)^2+(y_2-y_1)^2+(z_2-z_1)^2}$\\[3pt]
\tsl{ベクトル}$\vec a=(a_x,\ a_y,\ a_z)=a_x\vec i+a_y\vec j+a_z\vec k$\\[3pt]
\tsl{大きさ}$|\vec a|=\sqrt{a_x^2+a_y^2+a_z^2}$\\[3pt]
\tsl{内積}$\vec a\cdot\vec b=|\vec a||\vec b|\cos\theta=a_xb_x+a_yb_y+a_zb_z$\\[3pt]
\tsl{単位ベクトル}$\vec a\div|\vec a|=(\cos\theta_x,\ \cos\theta_y,\ \cos\theta_z)$（方向余弦）\\[3pt]
\tsl{平面}法線ベクトル $(a,\ b,\ c)$ の平面は $ax+by+cz=d$
\end{taisetsu}

\begin{kotoba}
\small
\bkword{次元}{場所を決めるのに必要な数の個数。直線は1、平面は2、空間は3。}
\bkword{右手系}{右手の親指・人さし指・中指の向きに、$x$・$y$・$z$ 軸を並べたもの。}
\bkword{基本ベクトル}{軸の＋の向きの単位ベクトル $\vec i$、$\vec j$、$\vec k$。}
\bkword{方向余弦}{ベクトルと3本の軸のなす角の cos。単位ベクトルの成分。}
\bkword{ねじれの位置}{空間の2直線が、交わらず平行でもない関係。}
\bkword{法線ベクトル}{平面に垂直なベクトル。平面の向きを表す。}
\bkword{投影図}{真正面・真上・真横から見た図。1つの座標を消した図。}
\end{kotoba}

\bkkigocards{\bkkigocard{$\symbf a$}{太字のエー}{ベクトル。$\vec a$ と同じ。細い字の $a$ は大きさ}{boldface vector}{4-13}}
{\bkkigocard{$\vec i,\ \vec j,\ \vec k$}{アイ・ジェイ・ケイ}{$x$・$y$・$z$ 軸の向きの基本ベクトル}{unit vectors}{4-13}}

%% ---------------- まちがえやすいところ（図の読み方と軸の並べ方。距離と方向余弦のまちがいは、例題のあとに置く） ----------------
\bksection{tsumazuki}{まちがえやすいところをチェック}

\begin{tsumazuki}{ななめにかいた $x$ 軸の長さを、目で読む}
\bkng{図で $x$ 方向の辺が短く見えるので、OA＝2 と読む}
\bkok{長さは、図ではなく座標で読む。A$(4,\ 0,\ 0)$ なら OA＝4}
見取図では、手前に出る $x$ 軸を、ななめに短くかいています。図の上の長さをものさしではかっても、本当の長さにはなりません。
\end{tsumazuki}

\begin{tsumazuki}{軸を左手系に並べてしまう}
\bkng{右を $x$、手前を $y$、上を $z$ にする}
\bkok{右を $x$ にしたいなら、手前を $-y$（おくを $y$）にする}
右を $x$、手前を $y$ にすると、鏡にうつした左手系になります。親指・人さし指・中指で、いつもたしかめましょう。
\end{tsumazuki}

%% ---------------- 例題 ----------------
\bksection{reidai}{例題でたしかめよう}

\begin{reidai}[直方体の頂点の座標]
図の直方体は、Oを原点とし、辺が3本の軸にそっています。OA＝4、OC＝3、OD＝2 です。\par
（1）頂点B、F、Gの座標を答えましょう。\par
（2）点Fの真下の床（$xy$ 平面）の点の座標を答えましょう。\par
\noindent\begin{minipage}[c]{0.4\linewidth}
\centering
\begin{tikzpicture}[ivn3d=6.5mm]
  \bkivnaxes{5.2}{4.2}{2.9}
  \bkivnbox{4}{3}{2}{bkInk}
  \foreach \p/\n/\a in {{(0,0,0)}/O/north east, {(4,0,0)}/A/north east, {(4,3,0)}/B/north west, {(0,3,0)}/C/south east, {(0,0,2)}/D/east, {(4,0,2)}/E/east, {(4,3,2)}/F/west, {(0,3,2)}/G/south west}
    {\fill[bkInk] \p circle (0.6mm); \node[font=\sffamily\footnotesize, anchor=\a, inner sep=1.5pt] at \p {\n};}
\end{tikzpicture}
\end{minipage}\hfill
\begin{minipage}[c]{0.57\linewidth}
\kangaekata Oから、$x$（手前）・$y$（右）・$z$（上）の順に、それぞれ何進むかを読む（①）。\par
\shiki（1）B：手前へ4、右へ3、上へ0\par
\hspace*{4.6\zw}F：手前へ4、右へ3、上へ2\par
\hspace*{4.6\zw}G：手前へ0、右へ3、上へ2\par
\hspace*{2.6\zw}（2）真下の点は、$x$、$y$ が同じで、$z=0$\par
\kotae （1）B$(4,\ 3,\ 0)$、F$(4,\ 3,\ 2)$、G$(0,\ 3,\ 2)$\quad（2）$(4,\ 3,\ 0)$（点B）
\end{minipage}
\end{reidai}

\begin{jibun}
例題1の図で、OA＝5、OC＝2、OD＝3 の直方体にかえます。頂点B、E、Gの座標と、点Gの真下の床の点の座標を答えましょう。
\bkjans{B$(5,\ 2,\ 0)$、E$(5,\ 0,\ 3)$、G$(0,\ 2,\ 3)$。Gの真下は $(0,\ 2,\ 0)$（点C）}
\end{jibun}

\begin{reidai}[空間の2点の距離]
（1）原点Oと点P$(2,\ 3,\ 6)$ の距離を求めましょう。\par
（2）2点A$(1,\ -2,\ 3)$、B$(5,\ 0,\ 7)$ の距離を求めましょう。\par
\kangaekata 座標の差（あと－まえ）を2乗してたし、$\sqrt{\ }$ をとる（③）。\par
\shiki（1）$\mathrm{OP}=\sqrt{2^2+3^2+6^2}=\sqrt{4+9+36}=\sqrt{49}=7$\par
\hspace*{2.6\zw}（2）差は\enspace $5-1=4$、\enspace $0-(-2)=2$、\enspace $7-3=4$\par
\hspace*{4.6\zw}$\mathrm{AB}=\sqrt{4^2+2^2+4^2}=\sqrt{16+4+16}=\sqrt{36}=6$\par
\kotae （1）7\quad（2）6
\end{reidai}
（2）の $0-(-2)$ は、マイナスをひくので $0+2=2$ です（2-2）。

\begin{jibun}
（1）原点Oと点P$(1,\ 4,\ 8)$ の距離を求めましょう。\par
（2）2点A$(2,\ 1,\ -1)$、B$(4,\ -2,\ 5)$ の距離を求めましょう。
\bkjans{（1）\nobreak $\sqrt{1+16+64}=\sqrt{81}=9$\quad（2）\nobreak 差は2、$-3$、6。$\sqrt{4+9+36}=7$}
\end{jibun}

\begin{tsumazuki}{距離の式で、$z$ をわすれる}
\bkng{A$(1,\ 2,\ 3)$、B$(4,\ 6,\ 15)$ の距離を $\sqrt{3^2+4^2}=5$}
\bkok{$\sqrt{3^2+4^2+12^2}=\sqrt{169}=13$}
平面の式のくせで、高さの差を入れわすれます。空間の距離は、3つの差を全部2乗してたします。
\end{tsumazuki}

\begin{reidai}[3次元のベクトルの計算]
（1）2点A$(1,\ 0,\ 2)$、B$(3,\ 2,\ 3)$ について、$\overrightarrow{\mathrm{AB}}$ の成分と大きさを求めましょう。\par
（2）$\vec a=(1,\ 2,\ 2)$、$\vec b=(3,\ 0,\ -4)$ のとき、$\vec a+\vec b$、$2\vec a-\vec b$ と、$|\vec a|$、$|\vec b|$、$|2\vec a-\vec b|$ を求めましょう。\par
\kangaekata 2点を結ぶベクトルは「B－A」。計算は成分ごと（④）。大きさは成分を2乗してたして $\sqrt{\ }$。\par
\shiki（1）$\overrightarrow{\mathrm{AB}}=(3-1,\ 2-0,\ 3-2)=(2,\ 2,\ 1)$\par
\hspace*{4.6\zw}$|\overrightarrow{\mathrm{AB}}|=\sqrt{4+4+1}=\sqrt9=3$\par
\hspace*{2.6\zw}（2）$\vec a+\vec b=(1+3,\ 2+0,\ 2+(-4))=(4,\ 2,\ -2)$\par
\hspace*{4.6\zw}$2\vec a=(2,\ 4,\ 4)$ なので\par
\hspace*{4.6\zw}$2\vec a-\vec b=(2-3,\ 4-0,\ 4-(-4))=(-1,\ 4,\ 8)$\par
\hspace*{4.6\zw}$|\vec a|=\sqrt{1+4+4}=3$、\enspace $|\vec b|=\sqrt{9+0+16}=5$\par
\hspace*{4.6\zw}$|2\vec a-\vec b|=\sqrt{1+16+64}=\sqrt{81}=9$\par
\kotae （1）$(2,\ 2,\ 1)$、大きさ3\quad（2）$(4,\ 2,\ -2)$、$(-1,\ 4,\ 8)$、3、5、9
\end{reidai}

\begin{jibun}
（1）2点A$(2,\ 1,\ 0)$、B$(4,\ -1,\ 1)$ について、$\overrightarrow{\mathrm{AB}}$ の成分と大きさを求めましょう。\par
（2）$\vec a=(2,\ 2,\ 1)$、$\vec b=(0,\ -4,\ 3)$ のとき、$\vec a+\vec b$、$2\vec a-\vec b$ と、$|\vec a|$、$|\vec b|$、$|2\vec a-\vec b|$ を求めましょう。
\bkjans{（1）\nobreak $(2,\ -2,\ 1)$、大きさ3。（2）\nobreak $(2,\ -2,\ 4)$、$(4,\ 8,\ -1)$、$|\vec a|=3$、$|\vec b|=5$、$|2\vec a-\vec b|=\sqrt{16+64+1}=9$}
\end{jibun}

\begin{reidai}[内積となす角・直交]
（1）$\vec a=(1,\ 1,\ 0)$、$\vec b=(0,\ 1,\ 1)$ のなす角を求めましょう。\par
（2）1辺1の立方体で、対角線 $\vec c=(1,\ 1,\ 1)$ と、辺 $\vec i=(1,\ 0,\ 0)$ のなす角を、電卓で求めましょう。\par
（3）$\vec p=(1,\ 2,\ 3)$ と $\vec q=(3,\ 0,\ -1)$ は直交するか、調べましょう。\par
\kangaekata $\cos\theta=\dfrac{\vec a\cdot\vec b}{|\vec a||\vec b|}$（⑤）。内積は、$x$ どうし・$y$ どうし・$z$ どうしをかけてたす。\par
\shiki（1）$\vec a\cdot\vec b=1\times0+1\times1+0\times1=1$、\enspace $|\vec a||\vec b|=\sqrt2\times\sqrt2=2$\par
\hspace*{4.6\zw}$\cos\theta=\dfrac{1}{2}$ なので\enspace $\theta=60°$\par
\hspace*{2.6\zw}（2）$\vec c\cdot\vec i=1\times1+1\times0+1\times0=1$、\enspace $|\vec c|=\sqrt3$、\enspace $|\vec i|=1$\par
\hspace*{4.6\zw}$\cos\theta=\dfrac{1}{\sqrt3}\approx0.577$。電卓（DEG）で $\cos^{-1}0.577\approx54.7°$\par
\hspace*{2.6\zw}（3）$\vec p\cdot\vec q=1\times3+2\times0+3\times(-1)=3+0-3=0$\par
\kotae （1）60°\quad（2）約54.7°\quad（3）内積が0なので、直交する
\end{reidai}

\begin{jibun}
（1）$\vec a=(2,\ 1,\ -1)$、$\vec b=(1,\ 2,\ 1)$ のなす角を求めましょう。\par
（2）$\vec a=(1,\ 2,\ 2)$、$\vec b=(2,\ -1,\ 2)$ のなす角を、電卓で求めましょう。\par
（3）$\vec p=(2,\ -1,\ 3)$ と $\vec q=(1,\ 5,\ 1)$ は直交するか、調べましょう。
\bkjans{（1）\nobreak 内積 $2+2-1=3$、大きさはどちらも $\sqrt6$。$\cos\theta=\dfrac{3}{6}=\dfrac{1}{2}$ で60°。（2）\nobreak 内積 $2-2+4=4$、大きさはどちらも3。$\cos\theta=\dfrac{4}{9}\approx0.444$ で約63.6°。（3）\nobreak $2-5+3=0$ なので直交する}
\end{jibun}

\begin{reidai}[空間の仕事]
ドローンが荷物をつるして運びます。ドローンが荷物を引く力は $\vec F=(4,\ 3,\ 20)$（N）で、そのあいだに荷物は $\vec s=(10,\ 0,\ 5)$（m）動きました。\par
（1）ドローンの力が荷物にした仕事を求めましょう。\par
（2）荷物にはたらく重力は、真下向きに\qty{20}{N}です。重力がした仕事を求めましょう。\par
\kangaekata 仕事は、力と動いたベクトルの内積 $W=F_xs_x+F_ys_y+F_zs_z$（⑤、4-12）。真下向きは $z$ 成分がマイナス。\par
\shiki（1）$W=4\times10+3\times0+20\times5=40+0+100=140$\par
\hspace*{2.6\zw}（2）重力は $(0,\ 0,\ -20)$ なので\par
\hspace*{4.6\zw}$W=0\times10+0\times0+(-20)\times5=-100$\par
\kotae （1）\qty{140}{J}\quad（2）\qty{-100}{J}
\end{reidai}
（1）で、$y$ 向きの力3\,Nは、荷物が $y$ の向きに動いていないので、仕事をしません。（2）の重力は、荷物が\qty{5}{m}上がるあいだ、動きと反対向きにはたらくので、仕事がマイナスです（4-12）。

\begin{jibun}
荷物を引く力が $\vec F=(5,\ -2,\ 30)$（N）で、荷物が $\vec s=(6,\ 4,\ 2)$（m）動きました。力がした仕事を求めましょう。また、真下向き\qty{30}{N}の重力がした仕事を求めましょう。
\bkjans{$5\times6+(-2)\times4+30\times2=30-8+60=\qty{82}{J}$。重力は $(-30)\times2=\qty{-60}{J}$}
\end{jibun}

\begin{reidai}[単位ベクトルと方向余弦]
$\vec a=(2,\ 3,\ 6)$ と同じ向きの単位ベクトル $\vec e$ を求めましょう。また、$\vec a$ と $z$ 軸のなす角 $\theta_z$ を、電卓で求めましょう。\par
\kangaekata $\vec a$ をその大きさでわる（⑥）。$\vec e$ の $z$ 成分が $\cos\theta_z$。\par
\shiki $|\vec a|=\sqrt{2^2+3^2+6^2}=\sqrt{49}=7$\par
\hspace*{2.6\zw}$\vec e=\dfrac{\vec a}{7}=\Bigl(\dfrac{2}{7},\ \dfrac{3}{7},\ \dfrac{6}{7}\Bigr)$\par
\hspace*{2.6\zw}$\cos\theta_z=\dfrac{6}{7}\approx0.857$。電卓（DEG）で $\cos^{-1}0.857\approx31.0°$\par
\kotae $\vec e=\Bigl(\dfrac{2}{7},\ \dfrac{3}{7},\ \dfrac{6}{7}\Bigr)$、\enspace $\theta_z\approx31.0°$
\end{reidai}
たしかめ：$\vec e$ の大きさは $\sqrt{\dfrac{4}{49}+\dfrac{9}{49}+\dfrac{36}{49}}=\sqrt{\dfrac{49}{49}}=1$ です。

\begin{jibun}
$\vec a=(1,\ 4,\ 8)$ と同じ向きの単位ベクトルを求めましょう。また、$\vec a$ と $z$ 軸のなす角を、電卓で求めましょう。
\bkjans{$|\vec a|=9$。$\Bigl(\dfrac{1}{9},\ \dfrac{4}{9},\ \dfrac{8}{9}\Bigr)$。$\cos\theta_z=\dfrac{8}{9}\approx0.889$ で約27.3°}
\end{jibun}

\begin{tsumazuki}{方向余弦を、そのまま角度だと思う}
\bkng{$\vec e=\Bigl(\dfrac{2}{7},\ \dfrac{3}{7},\ \dfrac{6}{7}\Bigr)$ だから、$z$ 軸とのなす角は $\dfrac{6}{7}$ 度}
\bkok{$\cos\theta_z=\dfrac{6}{7}$ なので、$\theta_z=\cos^{-1}\dfrac{6}{7}\approx31.0°$}
単位ベクトルの成分は、角の cos の値です。角度にするには、電卓の $\cos^{-1}$ を使います（DEG のモードで）。
\end{tsumazuki}

\begin{reidai}[法線ベクトルと平面の式]
点P$_0(1,\ 1,\ 1)$ を通り、$\vec n=(1,\ 2,\ 2)$ に垂直な平面の式を求めましょう。また、点Q$(3,\ 0,\ 1)$ はこの平面の上にあるか、調べましょう。\par
\kangaekata 平面の上の点を P$(x,\ y,\ z)$ とすると、$\vec n\cdot\overrightarrow{\mathrm{P_0P}}=0$（⑦）。\par
\shiki $\overrightarrow{\mathrm{P_0P}}=(x-1,\ y-1,\ z-1)$\par
\hspace*{2.6\zw}$1\times(x-1)+2(y-1)+2(z-1)=0$\par
\hspace*{2.6\zw}$x-1+2y-2+2z-2=0$\par
\hspace*{2.6\zw}$x+2y+2z=5$\par
\hspace*{2.6\zw}Qを入れると\enspace $3+2\times0+2\times1=5$\enspace で、式をみたす\par
\kotae $x+2y+2z=5$。Qは平面の上にある
\end{reidai}

\begin{jibun}
点P$_0(2,\ 0,\ 1)$ を通り、$\vec n=(3,\ 1,\ 2)$ に垂直な平面の式を求めましょう。また、点Q$(0,\ 2,\ 3)$ はこの平面の上にあるか、調べましょう。
\bkjans{$3(x-2)+y+2(z-1)=0$ より $3x+y+2z=8$。Qは $0+2+6=8$ で、平面の上にある}
\end{jibun}

%% ---------------- 物理・工学 ----------------
\begin{butsuri}
\noindent\begin{minipage}[c]{0.66\linewidth}
\textbf{半導体の結晶}\par\noindent コンピューターのチップに使うシリコンの結晶では、原子が立方体のくり返しの形に並んでいます。立方体の1辺は\qty{0.543}{nm}（ナノメートル。\qty{1}{nm}は10億分の1m）。1辺を1としたとき、角の原子 $(0,\ 0,\ 0)$ のとなりの原子は $\bigl(\frac{1}{4},\ \frac{1}{4},\ \frac{1}{4}\bigr)$ の所にあります。その距離は
\[ \sqrt{3}\times\frac{1}{4}\times\qty{0.543}{nm}\approx\qty{0.235}{nm} \]
で、実験ではかったシリコン原子どうしの距離と一致します。
\end{minipage}\hfill
\begin{minipage}[c]{0.31\linewidth}
\centering
\begin{tikzpicture}[ivn3d=24mm]
  \bkivnbox{1}{1}{1}{bkGray!70}
  \draw[bkPurple, line width=1.4pt] (0,0,0) -- (0.25,0.25,0.25);
  \shade[ball color=bkBlue!60] (0,0,0) circle (1.5mm);
  \shade[ball color=bkOrange!70] (0.25,0.25,0.25) circle (1.5mm);
  \node[font=\sffamily\footnotesize, anchor=north east, inner sep=1pt] at (0,0,0) {原点};
\end{tikzpicture}
\end{minipage}
\end{butsuri}

\begin{butsuri}
\noindent\begin{minipage}[c]{0.66\linewidth}
\textbf{3Dプリンター}\par\noindent とかした樹脂を出すノズルを、$(x,\ y,\ z)$ の数の指示どおりに動かします。たとえば「\mbox{X10 Y20 Z0.2}」は、「$x$ が\qty{10}{mm}、$y$ が\qty{20}{mm}、高さ\qty{0.2}{mm}の所へ」という意味です。高さ\qty{0.2}{mm}ずつ層を重ねて、立体を作ります。
\end{minipage}\hfill
\begin{minipage}[c]{0.31\linewidth}
\centering
\begin{tikzpicture}[x=1mm, y=1mm]
  \draw[bkGray!80!black, line width=1pt] (0,0) -- (0,18) -- (32,18) -- (32,0);
  \draw[bkGray!80!black, line width=0.8pt] (0,12) -- (32,12);
  \fill[bkOrange!70!black] (14,9) rectangle (19,13);
  \fill[bkOrange!70!black] (15.5,7) -- (17.5,7) -- (16.5,5.5) -- cycle;
  \fill[bkBlue!40] (8,0) rectangle (25,2);
  \fill[bkBlue!70] (11,2) rectangle (22,4.2);
  \node[font=\sffamily\footnotesize, anchor=west, inner sep=1pt] at (19.5,11) {ノズル};
\end{tikzpicture}
\end{minipage}
\end{butsuri}

\begin{butsuri}
\noindent\begin{minipage}[c]{0.66\linewidth}
\textbf{GPS（衛星測位）}\par\noindent 高さ約\qty{20000}{km}を回るGPS衛星は、自分の位置と時刻を電波で送っています。スマートフォンは、電波がとどくまでの時間から、衛星までの距離 $d_1$ を計算します。衛星の位置を $(x_1,\ y_1,\ z_1)$ とすると、自分の位置 $(x,\ y,\ z)$ は、空間の距離の式（③）で
\[ (x-x_1)^2+(y-y_1)^2+(z-z_1)^2=d_1^2 \]
をみたします。わからない数は $x$、$y$、$z$ と、スマートフォンの時計のずれの4つ。だから4つ以上の衛星の電波を使います。
\end{minipage}\hfill
\begin{minipage}[c]{0.31\linewidth}
\centering
\begin{tikzpicture}[x=1mm, y=1mm]
  \fill[bkBlue!15] (16,0) circle (9);
  \draw[bkBlue!70!black, line width=0.6pt] (16,0) circle (9);
  \fill[bkOrange] (16,9) circle (1);
  \foreach \sx/\sy in {2/30, 16/34, 30/30} {
    \draw[bkGray, line width=0.4pt, dashed] (16,9) -- (\sx,\sy);
    \fill[bkInk!70] (\sx-1.2,\sy-1) rectangle (\sx+1.2,\sy+1);
    \fill[bkBlue!60] (\sx-4,\sy-0.6) rectangle (\sx-1.6,\sy+0.6);
    \fill[bkBlue!60] (\sx+1.6,\sy-0.6) rectangle (\sx+4,\sy+0.6);}
  \node[font=\sffamily\footnotesize, anchor=north] at (16,-9.5) {地球};
\end{tikzpicture}
\end{minipage}
\end{butsuri}

\begin{column}{ドローンショーは、3つの数の設計図}
\noindent\begin{minipage}[c]{0.66\linewidth}
夜空に何百機ものドローンを光らせて絵をえがくドローンショーでは、絵の点の1つ1つに座標 $(x,\ y,\ z)$ をわりあてます。大切なのは、ドローンどうしがぶつからないことです。たとえば、2機の行き先が $(10,\ 5,\ 30)$ と $(11,\ 6,\ 31)$（単位はm）なら、その距離は③の式で
\[ \sqrt{1^2+1^2+1^2}=\sqrt3\approx1.73 \]
です。「どの2機も\qty{2}{m}以上はなす」と決めていたら、近すぎるので行き先を変えます。すべての組について、この計算をくり返して確かめます。
\end{minipage}\hfill
\begin{minipage}[c]{0.3\linewidth}
\centering
\begin{tikzpicture}[x=1mm, y=1mm]
  \fill[bkNavy!85] (-17,-17) rectangle (17,13);
  \foreach \t in {0,20,...,340} {\fill[bkYellow] ({16*sin(\t)^3*0.8},{(13*cos(\t)-5*cos(2*\t)-2*cos(3*\t)-cos(4*\t))*0.8}) circle (0.9);}
\end{tikzpicture}
\end{minipage}
\end{column}

%% ---------------- 章末問題 ----------------
\bkendproblems
\bkproblemgroup{基本}{3}

\bkproblem 例題1の図で、OA＝3、OC＝5、OD＝4 の直方体にかえます。頂点F、Gの座標と、点Fの真下の床の点の座標を答えましょう（\bkreidairef{1}）。
\bkpans{F$(3,\ 5,\ 4)$、G$(0,\ 5,\ 4)$。Fの真下は $(3,\ 5,\ 0)$}

\bkproblem 次の距離を求めましょう（\bkreidairef{2}）。\\
（1）原点と点P$(4,\ 4,\ 7)$\qquad（2）2点A$(-1,\ 2,\ 0)$、B$(1,\ 5,\ 6)$
\bkpans{（1）\nobreak $\sqrt{16+16+49}=9$\quad（2）\nobreak $\sqrt{4+9+36}=7$}

\bkproblem $\vec a=(1,\ -2,\ 2)$、$\vec b=(3,\ 4,\ 0)$ のとき、$\vec a+\vec b$、$2\vec a-\vec b$ と、$|\vec a|$、$|\vec b|$、$|2\vec a-\vec b|$ を求めましょう（\bkreidairef{3}）。
\bkpans{$(4,\ 2,\ 2)$、$(-1,\ -8,\ 4)$、$|\vec a|=3$、$|\vec b|=5$、$|2\vec a-\vec b|=\sqrt{1+64+16}=9$}

\bkproblem （1）$\vec a=(1,\ -1,\ 0)$、$\vec b=(1,\ 0,\ -1)$ のなす角を求めましょう。\\
（2）$\vec p=(4,\ 1,\ -2)$ と $\vec q=(1,\ 2,\ 3)$ は直交するか、調べましょう（\bkreidairef{4}）。
\bkpans{（1）\nobreak 内積1、大きさの積2。\mbox{$\cos\theta=\dfrac{1}{2}$ で60°}。（2）\nobreak $4+2-6=0$ なので直交する}

\bkproblem $\vec a=(2,\ -6,\ 3)$ と同じ向きの単位ベクトルを求めましょう。また、$\vec a$ と $x$ 軸のなす角を、電卓で求めましょう（\bkreidairef{6}）。
\bkpans{$|\vec a|=7$。$\Bigl(\dfrac{2}{7},\ -\dfrac{6}{7},\ \dfrac{3}{7}\Bigr)$。$\cos\theta_x=\dfrac{2}{7}\approx0.286$ で約73.4°}

\bkproblem 点P$_0(1,\ 2,\ 3)$ を通り、$\vec n=(2,\ -1,\ 1)$ に垂直な平面の式を求めましょう。また、点Q$(0,\ 0,\ 3)$ はこの平面の上にあるか、調べましょう（\bkreidairef{7}）。
\bkpans{$2(x-1)-(y-2)+(z-3)=0$ より $2x-y+z=3$。Qは $0-0+3=3$ で、平面の上にある}

\bkproblemgroup{チャレンジ}{4}

\bkproblem 校庭の点Oから飛び立ったドローンが、東へ\qty{30}{m}、北へ\qty{40}{m}、上へ\qty{120}{m}の点Bまで、まっすぐ飛びました。東・北・上を $x$・$y$・$z$ の＋の向きとします。（1）OBの距離を求めましょう。（2）$\overrightarrow{\mathrm{OB}}$ と真上（$z$ 軸）のなす角を、電卓で求めましょう（\bkreidairef{2}・\bkreidairef{6}）。
\bkpans{（1）\nobreak $\sqrt{900+1600+\num{14400}}=\sqrt{\num{16900}}=\qty{130}{m}$。（2）\nobreak $\cos\theta_z=\dfrac{120}{130}\approx0.923$ で約22.6°}

\bkproblem 荷物を引く力が $\vec F=(3,\ -4,\ 25)$（N）で、荷物が $\vec s=(20,\ 10,\ 4)$（m）動きました。力がした仕事と、真下向き\qty{25}{N}の重力がした仕事を求めましょう（\bkreidairef{5}）。
\bkpans{$60-40+100=\qty{120}{J}$。重力は $(-25)\times4=\qty{-100}{J}$}

\bkproblem 1辺1の立方体で、原点から向かいの角へ引いた対角線 $\vec a=(1,\ 1,\ 1)$ と、別の対角線 $\vec b=(-1,\ 1,\ 1)$ のなす角を、電卓で求めましょう（\bkreidairef{4}）。
\bkpans{内積 $-1+1+1=1$、大きさはどちらも $\sqrt3$。$\cos\theta=\dfrac{1}{3}\approx0.333$ で約70.5°}

\bkproblem 2点A$(1,\ 1,\ 0)$、B$(3,\ -1,\ 1)$ について、$\overrightarrow{\mathrm{AB}}$ の成分と大きさ、$\overrightarrow{\mathrm{AB}}$ と同じ向きの単位ベクトルを求めましょう（\bkreidairef{3}・\bkreidairef{6}）。
\bkpans{$(2,\ -2,\ 1)$、大きさ3。単位ベクトル $\Bigl(\dfrac{2}{3},\ -\dfrac{2}{3},\ \dfrac{1}{3}\Bigr)$}

\bkproblem 3Dプリンターのノズルが、点 $(10,\ 20,\ 0.2)$ から点 $(40,\ 60,\ 0.2)$ まで、まっすぐ動きました（単位はmm）。動いた距離を求めましょう（\bkreidairef{2}）。
\bkpans{差は30、40、0。$\sqrt{900+1600+0}=\sqrt{2500}=\qty{50}{mm}$}

\bkproblem $\vec a=(2,\ 2,\ 1)$ と $\vec b=(1,\ 2,\ 2)$ のなす角を、電卓で求めましょう。また、$\vec c=(1,\ -1,\ 0)$ と $\vec a$ は直交するか、調べましょう（\bkreidairef{4}）。
\bkpans{内積 $2+4+2=8$、大きさはどちらも3。$\cos\theta=\dfrac{8}{9}\approx0.889$ で約27.3°。$\vec c\cdot\vec a=2-2+0=0$ なので直交する}

\begin{column}{VRゴーグルは、6つの数をはかっている}
VR（バーチャルリアリティー）のゴーグルをつけて頭を動かすと、画面の景色もいっしょに動きます。ゴーグルは、頭の場所 $(x,\ y,\ z)$ の3つの数に加えて、頭がどちらを向いているか（上下・左右へのかたむきと、首のひねり）の3つの角も、1秒に何十回以上もはかっています。場所の3つと向きの3つで、合わせて6つの数です。ゲーム機のコントローラーや、ロボットの手先の動きも、同じように6つの数で表します。この「向き」の表し方は、第7巻の行列で学びます。
\begin{bkfigure}
\begin{tikzpicture}[x=1mm, y=1mm]
  \filldraw[fill=bkGray!25, draw=bkInk!70, line width=0.7pt, rounded corners=3mm] (-16,-6) rectangle (16,6);
  \fill[bkNavy!70] (-12,-3.5) rectangle (-2,3.5);
  \fill[bkNavy!70] (2,-3.5) rectangle (12,3.5);
  \begin{scope}[shift={(-46,0)}]
    \draw[bk arrow, bkRed, line width=1.1pt] (0,0) -- (12,0);
    \draw[bk arrow, bkBlue, line width=1.1pt] (0,0) -- (0,12);
    \draw[bk arrow, bkGreen!70!black, line width=1.1pt] (0,0) -- (-7,-7);
    \node[font=\sffamily\footnotesize, anchor=north, align=center] at (2,-9) {場所の3つの数\\$(x,\ y,\ z)$};
  \end{scope}
  \begin{scope}[shift={(44,0)}]
    \draw[bkOrange!80!black, line width=1.1pt, -{Stealth[length=4pt]}] (-8,4) arc (160:20:8.5 and 3);
    \draw[bkOrange!80!black, line width=1.1pt, -{Stealth[length=4pt]}] (8,-7) arc (-60:60:3 and 8);
    \draw[bkOrange!80!black, line width=1.1pt, -{Stealth[length=4pt]}] (-6,-8) arc (220:320:8 and 3);
    \node[font=\sffamily\footnotesize, anchor=north, align=center] at (0,-11) {向きの3つの角};
  \end{scope}
\end{tikzpicture}
\end{bkfigure}
\end{column}

\begin{furikaeri}
\begin{checklist}
\item 空間の点を $(x,\ y,\ z)$ で表し、2点の距離を求められる。\bkfurref{bk:4-13:zahyo}
\item 3次元のベクトルを成分で計算し、内積からなす角を求められる。\bkfurref{bk:4-13:naiseki}
\item 基本ベクトル・単位ベクトル・法線ベクトルで、向きや平面を表せる。\bkfurref{bk:4-13:kihon}
\end{checklist}
\end{furikaeri}

\bkfukidashi{sora}{3つの数で、空間のどこでも言えるようになった。次は、ベクトルどうしの新しいかけ算「外積」だよ。}
